\section{AMC 10 10A 2024}
        \begin{enumerate}
        	\item (\textit{AMC 10 10A 2024}) What is the value of $9,901\cdot101-99\cdot10,101?$


 $\textbf{(A)}~2\qquad\textbf{(B)}~20\qquad\textbf{(C)}~200\qquad\textbf{(D)}~202\qquad\textbf{(E)}~2020$


 
 \textbf{Answer: }A



\textbf{Solution 1 (Direct Computation)}

The likely fastest method will be direct computation. $9901\cdot101$ evaluates to $1000001$ and $99\cdot10101$ evaluates to $999999$. The difference is $\boxed{\textbf{(A) }2}.$

 Solution by \href{https://artofproblemsolving.com/wiki/index.php?title=User:Juwushu}{juwushu}.

 



\textbf{Solution 2 (Distributive Property)}

We have
 \begin{align*} 9901\cdot101-99\cdot10101 &= (10000-99)\cdot101-99\cdot(10000+101) \\ &= 10000\cdot101-99\cdot101-99\cdot10000-99\cdot101 \\ &= (10000\cdot101-99\cdot10000)-2\cdot(99\cdot101) \\ &= 2\cdot10000-2\cdot9999 \\ &= \boxed{\textbf{(A) }2}. \end{align*}
~MRENTHUSIASM

 



\textbf{Solution 3 (Solution 1 but Distributive)}

Note that $9901\cdot101=9901\cdot100+9901=990100+9901=1000001$ and  $99\cdot10101=100\cdot10101-10101=1010100-10101=999999$, therefore the answer is $1000001-999999=\boxed{\textbf{(A) }2}$.

 



\textbf{Solution 4 (Modular Arithmetic)}

Evaluating the given expression $\pmod{10}$ yields $1-9\equiv 2 \pmod{10}$, so the answer is either $\textbf{(A)}$ or $\textbf{(D)}$. Evaluating $\pmod{101}$ yields $0-99\equiv 2\pmod{101}$. Because answer $\textbf{(D)}$ is $202=2\cdot 101$, that cannot be the answer, so we choose choice $\boxed{\textbf{(A) }2}$.

 



\textbf{Solution 5 (Process of Elimination)}

We simply look at the units digit of the problem we have (or take mod $10$)
 \[9901\cdot101-99\cdot10101 \equiv 1\cdot1 - 9\cdot1 = 2 \mod{10}.\]Since the only answers with $2$ in the units digit are $\textbf{(A)}$ and $\textbf{(D)}$, we can then continue if you are desperate to use guess and check or an actually valid method to find the answer is $\boxed{\textbf{(A) }2}$.

 



\textbf{Solution 6 (Faster Distribution)}

Observe that $9901=9900+1=99\cdot100+1$ and $10101=10100+1=101\cdot100+1$
 \begin{align*} \Rightarrow9901\cdot101-99\cdot10101 & = ((9900\cdot101)+(1\cdot101))-((99\cdot10100)+(99\cdot1)) \\ &=(99\cdot100\cdot101)+101-(99\cdot100\cdot101)-99 \\ &=101-99 \\ &=\boxed{\textbf{(A) }2}. \end{align*}


 



\textbf{Solution 7 (Cubes)}

Let $x=100$. Then, we have\begin{align*}101\cdot 9901=(x+1)\cdot (x^2-x+1)=x^3+1, \\99\cdot 10101=(x-1)\cdot (x^2+x+1)=x^3-1.\end{align*}
Then, the answer can be rewritten as $(x^3+1)-(x^3-1)= \boxed{\textbf{(A) }2}.$

 



\textbf{Solution 8 (Super Fast)}

It's not hard to observe and express $9901$ into $99\cdot100+1$, and $10101$ into $101\cdot100+1$.

 We then simplify the original expression into $(99\cdot100+1)\cdot101-99\cdot(101\cdot100+1)$, which could then be simplified into $99\cdot100\cdot101+101-99\cdot100\cdot101-99$, which we can get the answer of $101-99=\boxed{\textbf{(A) }2}$.

 



\textbf{Solution 9 (Estimation) *VERY FAST*}

Notice that the answer choices are significantly different in value. This allows us to estimate the answer. $9901$ is about $10000$, and $101$ is about $100$. $99$ is about $100$, and $10101$ is about $10000$. Computing, we get $10000 \cdot 100-100 \cdot 10000 = 0$. The closest answer to this estimation is $\boxed{\textbf{(A) }2}$.

 



\textbf{Solution 10}

We can see that the units digit of the expression is $1 - 9 \Rightarrow 11 - 9 \Rightarrow 2$, elimination options B, C, and E. Next, notice that $(9901)(101)$ is divisible by 101 while $(99)(10101)$ is not divisible by 101 (to see this, notice that 101 is prime, and $10101 = 10100 + 1$, so is not divisible by 101). This means that the final answer is not divisible by 101, eliminating $\textbf{(D)}$, so the answer is $\boxed{\textbf{(A) }2}$.

 


	\item (\textit{AMC 10 10A 2024}) A model used to estimate the time it will take to hike to the top of the mountain on a trail is of the form $T=aL+bG,$ where $a$ and $b$ are constants, $T$ is the time in minutes, $L$ is the length of the trail in miles, and $G$ is the altitude gain in feet. The model estimates that it will take $69$ minutes to hike to the top if a trail is $1.5$ miles long and ascends $800$ feet, as well as if a trail is $1.2$ miles long and ascends $1100$ feet. How many minutes does the model estimates it will take to hike to the top if the trail is $4.2$ miles long and ascends $4000$ feet?


 $\textbf{(A) }240\qquad\textbf{(B) }246\qquad\textbf{(C) }252\qquad\textbf{(D) }258\qquad\textbf{(E) }264$


 
 \textbf{Answer: }B



\textbf{Solution 1}

Plug in the values into the equation to give you the following two equations:\begin{align*}69&=1.5a+800b, \\69&=1.2a+1100b.\end{align*}
Solving for the values $a$ and $b$ gives you that $a=30$ and $b=\frac{3}{100}$. These values can be plugged back in showing that these values are correct.Now, using the given length of the trail, $4.2$, and the given vertical increase, $4000$ , we get a final answer of $\boxed{\textbf{(B) }246}.$

 Solution by \href{https://artofproblemsolving.com/wiki/index.php?title=User:Juwushu}{juwushu}.

 Minor edits by ParticlePhysics and TigerSenju

 



\textbf{Solution 2}

Alternatively, observe that using $a=10x$ and $b=\frac{y}{100}$ makes the numbers much closer to each other in terms of magnitude.

 Plugging in the new variables:\begin{align*}69&=15x+8y, \\69&=12x+11y.\end{align*}


 The solution becomes more obvious in this way, with $15+8=12+11=23$, and since $23\cdot 3=69$, we determine that $x=y=3$.

 The question asks us for $4.2a+4000b=42x+40y$. Since $x=y$, we have $(40+42)\cdot 3=\boxed{\textbf{(B) }246}$.

 


	\item (\textit{AMC 10 10A 2024}) Let $n$ be the least prime number that can be written as the sum of $5$ distinct prime numbers. What is
the sum of the digits of $n$?


 $\textbf{(A) }5\qquad\textbf{(B) }7\qquad\textbf{(C) }8\qquad\textbf{(D) }10\qquad\textbf{(E) }13$


 
 \textbf{Answer: }B



\textbf{Solution 1}

Let the requested sum be $S.$ Recall that $2$ is the only even (and the smallest) prime, so $S$ is odd. It follows that the five distinct primes are all odd. The first few odd primes are $3,5,7,11,13,17,19,\ldots.$ We conclude that $S>3+5+7+11+13=39,$ as $39$ is a composite. The next possible value of $S$ is $3+5+7+11+17=43,$ which is a prime. Therefore, we have $S=43,$ and the sum of its digits is $4+3=\boxed{\textbf{(B) }7}.$

 



\textbf{Solution 2}

We notice that the minimum possible value of the sum of $5$ odd distinct primes is $3 + 5 + 7 + 11 + 13 = 39$, which is not a prime. The smallest prime greater than that is $41$. However, this cannot be written as the sum of $5$ distinct primes, since $15$ is not prime. However, $43$ can be written as $3 + 5 + 7 + 11 + 17 = 43$, so the answer is $4 + 3 = \boxed{\textbf{(B) }7}$ 

 


	\item (\textit{AMC 10 10A 2024}) The number $2024$ is written as the sum of not necessarily distinct two-digit numbers. What is the least number of two-digit numbers needed to write this sum?


 $\textbf{(A) }20\qquad\textbf{(B) }21\qquad\textbf{(C) }22\qquad\textbf{(D) }23\qquad\textbf{(E) }24$


 
 \textbf{Answer: }B



\textbf{Solution 1}

Since we want the least number of two-digit numbers, we maximize the two-digit numbers by choosing as many $99$s as possible. Since $2024=99\cdot20+44\cdot1,$ we choose twenty $99$s and one $44,$ for a total of $\boxed{\textbf{(B) }21}$ two-digit numbers.

 



\textbf{Solution 2}

We claim the answer is $21$. This can be achieved by adding twenty $99$'s and a $44$. To prove that the answer cannot be less than or equal to $20$, we note that the maximum value of the sum of $20$ or less two digit numbers is $20 \cdot 99 = 1980$, which is smaller than $2024$, so we are done. Thus, the answer is $\boxed{\textbf{(B) }21}$.

 



\textbf{Solution 3 (Same as Solution 1 but Using 100=99+1)}

$2024=100\cdot20+24$. Since $100=99+1$, $2024=(99+1)\cdot20+24=99\cdot20+1\cdot20+24=99\cdot20+44$. Therefore a total of $\boxed{\textbf{(B) }21}$ two-digit numbers are needed.

 



\textbf{Solution 4}

The maximum $2$-digit number is $99$, but try $100$. $2024 \div 100$ is a little more than $20$, and the remainder is less than $100$, by intuition, so there's $20 +$ the remainder $1 = \boxed{\textbf{(B) }21}$.

 


	\item (\textit{AMC 10 10A 2024}) What is the least value of $n$ such that $n!$ is a multiple of $2024$?


 $\textbf{(A) } 11\qquad\textbf{(B) } 21\qquad\textbf{(C) } 22\qquad\textbf{(D) } 23\qquad\textbf{(E) } 253$


 
 \textbf{Answer: }D



\textbf{Solution}

Note that $2024=2^3\cdot11\cdot23$ in the prime factorization. Since $23!$ is a multiple of $2^3, 11,$ and $23,$ we conclude that $23!$ is a multiple of $2024.$ Therefore, we have $n=\boxed{\textbf{(D) } 23}.$

 <u>\textbf{Remark}</u>

 Memorizing the prime factorization of the current year is useful for the AMC 8/10/12 Exams.

 <u>\textbf{Remark}</u>

 Note that $2024$ is $2025 -1$, which is $45^2 -1 = (45+1)(45-1) = 46\cdot44$.

 


	\item (\textit{AMC 10 10A 2024}) What is the minimum number of successive swaps of adjacent letters in the string $ABCDEF$ that are needed to change the string to $FEDCBA?$ (For example, $3$ swaps are required to change $ABC$ to $CBA;$ one such sequence of swaps is
$ABC\to BAC\to BCA\to CBA.$)


 $\textbf{(A)}~6\qquad\textbf{(B)}~10\qquad\textbf{(C)}~12\qquad\textbf{(D)}~15\qquad\textbf{(E)}~24$


 
 \textbf{Answer: }D



\textbf{Solution 1 (Analysis)}

Procedurally, it takes:

 
                \begin{itemize}
                    \item $5$ swaps for $A$ to move to the sixth spot, giving $BCDEFA.$
                \end{itemize}


                
                \begin{itemize}
                    \item $4$ swaps for $B$ to move to the fifth spot, giving $CDEFBA.$
                \end{itemize}


                
                \begin{itemize}
                    \item $3$ swaps for $C$ to move to the fourth spot, giving $DEFCBA.$
                \end{itemize}


                
                \begin{itemize}
                    \item $2$ swaps for $D$ to move to the third spot, giving $EFDCBA.$
                \end{itemize}


                
                \begin{itemize}
                    \item $1$ swap for $E$ to move to the second spot (so $F$ becomes the first spot), giving $FEDCBA.$
                \end{itemize}


                Together, the answer is $5+4+3+2+1=\boxed{\textbf{(D)}~15}.$

 



\textbf{Solution 2 (Recursive Approach)}

We can proceed by a recursive tactic on the number of letters in the string.

 Looking at the string $A$, there are $0$ moves needed to change it to the string $A$

 Then, there is $1$ move to change $AB$ to $BA$.

 Similarly, there is $3$ moves needed for three letters (said in the problem).

 There are $6$ moves needed to change $ABCD$ to $DCBA$.

 We see a pattern of $0,1,3,6,...$. We notice that the difference between consecutive terms is increasing by $1$, so in the same way, for $5$ letters, we would need $10$ moves, and for $6$, we would need $\boxed{\textbf{(D)}~15}$ moves.

 Thinking why, when we start making these moves, we see that for a string of length $n$, it takes $n-1$ moves to move the last letter to the front. After, we get a string that will be changed identically to a string of length $n-1$. This works in our pattern above and is another way to think about the problem!

 Note: 

 The sequence consists of triangular numbers shifted a term up (as it starts with $0$ on term $1$ and $1$ on term $2$.)Thus, our explicit formula is 
 \[\dfrac{(n-1)(n+1-1)}{2} = \dfrac{(n)(n-1)}{2}\] and as $n = 6$ in this case ($6$ letters), our answer is $\dfrac{(6)(6-1)}{2} = \boxed{\textbf{(D)}~15}$ 

 



\textbf{Solution 3 (Solution 2 Done Fast)}

We can see that the most efficient way to change $ABCDEF$ to $FEDCBA$ is the same as changing $ABCDE$ to $EDCBA$ and then moving $F$ to the front in $5$ moves. Similarly, this would carry on downwards, where to change $ABCDE$ to $EDCBA$ would be to change $ABCD$ to $DCBA$ and move $E$ $4$ times to the front. This pattern will carry on until $AB$ to $BA$ would be $1$, and $A$ to $A$ would be $0$. The answer would be $0(A)+1(B)+2(C)+3(D)+4(E)+5(F)$, which is $\boxed{\textbf{(D)}~15}$ moves.

 



\textbf{Solution 4 (If you don't notice anything)}

Simply, just blindly doing it, the most straightforward way is to shift F all the way back. From ABCDEF:

 $ABCDFE ->ABCFDE ->ABFCDE ->AFBCDE ->FABCDE$

 For E:

 $FABCED -> FABECD -> FAEBCD -> FEABCD$

 For D:

 $FEABDC -> FEADBC -> FEDABC$

 For C:

 $FEDACB -> FEDCAB$

 For B:

 FEDCBA

 That's it, you don't need to do it with A. Still $\boxed{\textbf{(D)}~15}$ moves.

 



\textbf{Solution 5 (Inversions of Permutations)}

Let the letters $A, B, C, D, E, F$ correspond to the numbers $1, 2, 3, 4, 5, 6$ respectively. We want to find the number of swaps required to transform the permutation $1 2 3 4 5 6$ into the permutation $6 5 4 3 2 1$. 

 Here, we let $1 2 3 4 5 6$ be the natural order. Then the total number of inversions of the permutation $6 5 4 3 2 1$ is $1+2+3+4+5=15$. Hence the answer is $\boxed{\textbf{(D)}~15}$

 


	\item (\textit{AMC 10 10A 2024}) The product of three integers is $60$. What is the least possible positive sum of the 
three integers? 


 $\textbf{(A) }2\qquad\textbf{(B) }3\qquad\textbf{(C) }5\qquad\textbf{(D) }6\qquad\textbf{(E) }13$


 
 \textbf{Answer: }B



\textbf{Solution 1}

We notice that the optimal solution involves two negative numbers and a positive number. Thus we may split $60$ into three factors and choose negativity. We notice that $10\cdot6\cdot1=10\cdot(-6)\cdot(-1)=60$, and trying other combinations does not yield lesser results so the answer is $10-6-1=\boxed{\textbf{(B) }3}$.

 



\textbf{Solution 2}

We have $abc = 60$. Let $a$ be positive, and let $b$ and $c$ be negative. Then we need $a > |b + c|$. If $a = 6$, then $|b + c|$ is at least $7$, so this doesn't work. If $a = 10$, then $(b,c) = (-6,-1)$ works, giving $10 - 7 = \boxed{\textbf{(B) }3}$

 



\textbf{Solution 3}

We can see that the most optimal solution would be $1$ positive integer and $2$ negative ones (as seen in solution 1). Let the three integers be $x$, $y$, and $z$, and let $x$ be positive and $y$ and $z$ be negative. If we want the optimal solution, we want the negative numbers to be as large as possible. so the answer should be $-60 \cdot -1 \cdot 1$, where $-60 - 1 + 1 = -60$... right?

 No! Our sum must be positive, so that would be invalid! We see that -$30, -20, -15$, and $-10$ are too less to allow the sum to be positive. For example, $-15 = z. -15xy=60$, so $xy = -4$. For $xy$ to be the most positive, we will have $4$ and $-1$. Yet, $-15+4-1$ is still negative. After $-10$, the next factor of $60$ would be $-6$. if $z$ = $-6$, $xy = -10$. This might be positive! Now, if we have $z = -6, y = -1$, and $x = 10, x + y + z = 3$. It cannot be smaller because $x = 5$ and $y = -2$ would result in $x + y + z$ being negative. Therefore, our answer would be $\boxed{\textbf{(B) }3}$.

 Tiny minor edits made by SuperVince1

 


	\item (\textit{AMC 10 10A 2024}) Amy, Bomani, Charlie, and Daria work in a chocolate factory. On Monday Amy, Bomani, and Charlie started working at $1:00 PM$ and were able to pack $4$, $3$, and $3$ packages, respectively, every $3$ minutes. At some later time, Daria joined the group, and Daria was able to pack $5$ packages every $4$ minutes. Together, they finished packing $450$ packages at exactly $2:45 PM$. At what time did Daria join the group?


 $\textbf{(A) }1:25\text{ PM}\qquad\textbf{(B) }1:35\text{ PM}\qquad\textbf{(C) }1:45\text{ PM}\qquad\textbf{(D) }1:55\text{ PM}\qquad\textbf{(E) }2:05\text{ PM}$


 
 \textbf{Answer: }A



\textbf{Solution 1}

Note that Amy, Bomani, and Charlie pack a total of $4+3+3=10$ packages every $3$ minutes.

 The total amount of time worked is $1$ hour and $45$ minutes, which when converted to minutes, is $105$ minutes. This means that since Amy, Bomani, and Charlie worked for the entire $105$ minutes, they in total packed $\dfrac{105}{3}\cdot10=350$ packages.

 Since $450$ packages were packed in total, then Daria must have packed $450-350=100$ packages in total, and since she packs at a rate of $5$ packages per $4$ minutes, then Daria worked for $\dfrac{100}{5}\cdot4=80$ minutes, therefore Daria joined $80$ minutes before $2:45$ PM, which was at $\boxed{\textbf{(A) }1:25\text{ PM}}$.

 



\textbf{Solution 2}

Let the time, in minutes, elapsed between $1:00$ and the time Daria joined the packaging be $x$. Since Amy packages $4$ packages every $3$ minutes, she packages $\frac{4}{3}$ packages per minute. Similarly, we can see that both Bomani and Charlie package $1$ package per minute, and Daria packages $\frac{5}{4}$ packages every minute.

 Before Daria arrives, we can write the total packages packaged as $x(\frac{4}{3} + 1 + 1) = x(\frac{10}{3})$. Since there are $105$ minutes between $1:00$ and $2:45$, Daria works with the other three for $105-x$ minutes, meaning for that time there are $(105-x)(\frac{4}{3} + 1 + 1 + \frac{5}{4}) = (105-x)(\frac{55}{12})$ packages packaged. 

 Adding the two, we get $x(\frac{10}{3}) + (105-x)(\frac{55}{12}) = 450$ (The total packaged in the entire time is $450$). Solving this equation, we get $x=25$, meaning Daria arrived $25$ minutes after $1:00$, meaning the answer is $\boxed{\textbf{(A) }1:25\text{ PM}}$.

 


	\item (\textit{AMC 10 10A 2024}) In how many ways can $6$ juniors and $6$ seniors form $3$ disjoint teams of $4$ people so 
that each team has $2$ juniors and $2$ seniors? 


 $\textbf{(A) }720\qquad\textbf{(B) }1350\qquad\textbf{(C) }2700\qquad\textbf{(D) }3280\qquad\textbf{(E) }8100$


 
 \textbf{Answer: }B



\textbf{Solution 1}

The number of ways in which we can choose the juniors for the team are ${6\choose2}{4\choose2}{2\choose2}=15\cdot6\cdot1=90$. Similarly, the number of ways to choose the seniors are the same, so the total is $90\cdot90=8100$. But we must divide the number of permutations of three teams, since the order in which the teams were chosen doesn't matter, which is $3!$. Thus, the answer is $\frac{8100}{3!}=\frac{8100}{6}=\boxed{\textbf{(B) }1350}$

 



\textbf{Solution 2}

Let's first suppose that the three teams are distinguishable. The number of ways to choose the first team of two juniors and two seniors is ${6\choose2}{6\choose2} = 15\cdot 15 = 225$. Now, only four juniors and four seniors are left to choose the second team. Thus, the second team can be formed in ${4\choose2}{4\choose2} = 6\cdot 6 = 36$ ways. There are now only two juniors and two seniors left, so the third team can only be formed in one way. Thus, the total number of ways in which we can choose two juniors and two seniors for three distinguishable teams is $225\cdot 36 = 8100$ ways. However, the problem does not require the teams to be distinguishable. Therefore, we must divide by $3!=6$ to find the answer for three indistinguishable teams.

 The answer is $\frac{8100}{6} =\boxed{\textbf{(B) }1350}$

 


	\item (\textit{AMC 10 10A 2024}) Consider the following operation. Given a positive integer $n$, if $n$ is a multiple of $3$, then you replace $n$ by $\frac{n}{3}$. If $n$ is not a multiple of $3$, then you replace $n$ by $n+10$. For example, beginning with $n=4$, this procedure gives $4\to14\to24\to8\to18\to6\to2\to12\to\cdots$. Suppose you start with $n=100$. What value results if you perform this operation exactly $100$ times?


 $\textbf{(A) }10\qquad\textbf{(B) }20\qquad\textbf{(C) }30\qquad\textbf{(D) }40\qquad\textbf{(E) }50$


 
 \textbf{Answer: }C



\textbf{Solution 1 (Fast Solution)}

Let $s$ be the number of times the operation is performed. Notice the sequence goes $100 \to 110 \to 120 \to 40 \to 50 \to 60 \to 20 \to 30 \to 10 \to 20 \to \cdots$. Thus, for $s \equiv 1 \pmod{3}$, the value is $30$. Since $100 \equiv 1 \pmod{3}$, the answer is $\boxed{\textbf{(C) }30}$.

 



\textbf{Solution 2 (More Explanatory)}

Looking at the first few values of our operation, we get $100 \to 110 \to 120 \to 40 \to 50 \to 60 \to 20 \to 30 \to 10 \to 20$. We can see that $30$ goes to $10$, then to $20$, then back to $30$, and the loop resets. After 7 operations, we reach $30$. We still have 93 operations left, so because the loop will run exactly $31$ times $(93/3)$, we will reach $30$ again. So, the answer is $\boxed{\textbf{(C) } 30}$.

 



\textbf{Solution 3 (very slightly different than previous)}

Calculating the first few values, we get $100 \to 110 \to 120 \to 40 \to 50 \to 60 \to 20 \to 30 \to 10 \to 20$. We can see that $20$ will go to $30$, then to $10$, then back to $20$, and then the loop resets. After $6$ moves, we reach $20$, the start of the cycle. We still have $100-5$ moves to go, so to find what number we land on after $95$ more steps, we can do  $95 \pmod {3} \equiv (9 + 4) \pmod {3} = 1$, meaning we go from $20 \to \boxed{\textbf{(C) } 30}$.

 


	\item (\textit{AMC 10 10A 2024}) How many ordered pairs of integers $(m, n)$ satisfy $\sqrt{n^2 - 49} = m$?


 $\textbf{(A)}~1\qquad\textbf{(B)}~2\qquad\textbf{(C)}~3\qquad\textbf{(D)}~4\qquad\textbf{(E)}~\text{infinitely many}$


 
 \textbf{Answer: }D



\textbf{Solution 1}

Note that $m$ is a nonnegative integer.

 We square, rearrange, and apply the difference of squares formula to the given equation: 
 \[(n+m)(n-m)=49.\]It is clear that $n+m\geq n-m,$ so $(n+m,n-m)=(49,1),(7,7),(-7,-7),(-1,-49).$ Each ordered pair $(n+m,n-m)$ gives one ordered pair $(m,n),$ so there are $\boxed{\textbf{(D)}~4}$ such ordered pairs $(m,n).$

 <u>\textbf{Remark}</u>

 From $(n+m,n-m)=(49,1),(7,7),(-7,-7),(-1,-49),$ we get $(m,n)=(24,25),(0,7),(0,-7),(24,-25),$ respectively.

 



\textbf{Solution 2}

Squaring both sides of the given equation gives 
 \[n^2-49=m^2\rightarrow n^2-m^2=49\rightarrow (n+m)(n-m)=49\] Splitting $49$ into its factors (keep in mind it doesn't ask for positive integers, so the factors can be double negative, too) gives six cases:

 $(1\cdot49)$

 $(7\cdot7)$

 $(49\cdot1)$

 $(-1\cdot -49)$

 $(-7\cdot -7)$

 $(-49\cdot -1)$.

 Note that the square root in the problem doesn't have $\pm$ with it. Therefore, if there are two solutions, $(n,m)$ and $(n,-m)$, then these together are to be counted as one solution.The solutions expressed as $(n,m)$ are:

 $(25,24)$

 $(25,-24)$

 $(7,0)$

 $(-7,0)$

 $(-25,24)$

 $(-25,-24)$.

 $(25,24)$ and $(25,-24)$ are to be counted as one, same for $(-25,24)$ and $(-25,-24)$. Therefore, the solution is $\boxed{\textbf{(D)}~4}$.

 



\textbf{Solution 3 (Crazy Rush)}

From looking at the problem, it's obvious that $(\pm7,0)$ are already solutions. From squaring and rearranging, we get 
 \[n^2-m^2=49.\] We know that the difference between two consecutive squares is always odd, and for each pair of increasing consecutive squares, the difference starts from $3$ and increases by $2$ each time. This means that there is an existing pair, $(\pm n,m)$ of consecutive squares that will satisfy this equation. 

 Also note that the answer cannot be infinity because the difference between two squares will increase as the two squares get higher, consecutive or not.

 Therefore, the solutions $(\pm7,0)$ and $(\pm n,m)$ where $n$ and $m$ are consecutive that have a square difference of $49,$ give the answer of $\boxed{\textbf{(D)}~4}.$

 Minor Note: Looking at the problem, it is obvious that $-7$ and $7$ are solutions and you can eliminate $1,3,$ and infinite. You can assume that there is more than $2$ solutions because when determining square expressions, these problems usually always have more than an obvious solution.

 



\textbf{Solution 4 (Quick)}

Clearly, $(7,0)$ and $(-7,0)$ are solutions.Notice that $49 = 7^2$. Remembering our Pythagorean triples, we realize $(25, 24)$ is a solution as well, and by extension, so is $(-25, 24)$. Since 7 is not part of any other Pythagorean triple, we can conclude the answer is  $\boxed{\textbf{(D)}~4}$.

 



\textbf{Solution 5 (idk why someone would ever use this)}

Since in $n^2-49$ there are no restrictions of the sign of $n$, we know that the solution will come in pairs, with negative and positive $n$. Therefore, the answer will not be $1$ or $3$. We can further confirm that it is not $1$, since the problem would probably ask for an operation of $n$ and $m$ instead of the number of solutions. We can also guess it is not $2$, since then the problem would ask for the sum of one of the variables in the solutions. It seems improbable that over the integers there would be infinite answers. So the answer is $\boxed{\textbf{(D) }4}$

 (do not ever use this)

 



\textbf{Solution 6 (Pythagorean Theorem)}

Rearranging the equation gives

 
 \[m^2+7^2=n^2\]

 This is a classic Pythagorean Theorem equation, where $7$ is a leg and $m$ and $n$ are the two remaining legs. All we have to do is to find triples that satisfy this equation. It is common knowledge that $(7, 24, 25)$ is a triple, thus $(24, 25)$ and $(24, -25)$ are two solutions. Another one is $(0, 7, 7)$, so $(0, 7)$ and $(0, -7)$ are two other solutions. We can be sure that there are no other triples to find because $4$ is the highest answer choice, otherwise it would be "infinitely many", which isn't plausible. Therefore, the answer is $\boxed{\textbf{(D) }4}$.

 -urAntimony (minor edit)

 



\textbf{Solution 7}

Square both sides giving $n^2-49=m^2$, then rearrange the equation into a quadratic: $n^2-49-m^2=0$. then find the Delta ($b^2-4ac$) where $a,b,$ and $c$ are the coefficients of the quadratic in the respective order. Calculating it we get $2401-4=2397$, since this is a positive number it means there are two solutions but both solutions can be both $m$ or $n$, so the answer is $\boxed{\textbf{(D) }4}$.

 



\textbf{Solution 8 (Graphing)}

We can rewrite the equation as $y=\sqrt{x^{2}-49}$ and graph it to find integer solutions. It is a hyperbola of standard form $\frac{x^{2}}{49}-\frac{y^{2}}{49}=1$. Note that the graph of the function does not exist below the y-axis. We find that where $\left|x\right|>25$ integer solutions are impossible, since the difference between the function and its asymptote $y=x$ in that region is less than one and always decreasing; thus, since $y=x$ is always an integer, $\sqrt{x^{2}-49}$ must be a fractional amount less than that integer, so we can rule out option E (infinite integer solutions). We notice that the roots are a solution $\left(\pm7,\ 0\right)$, as well as $\left(\pm25,\ 24\right)$, making four solutions. $4$ is already the maximum possible answer choice (having ruled out E), so we know there are no more solutions. Thus, the answer is $\boxed{\textbf{(D) }4}$.

 


	\item (\textit{AMC 10 10A 2024}) Zelda played the \textit{Adventures of Math} game on August 1 and scored $1,700$ points. She continued to play daily over the next $5$ days. The bar chart below shows the daily change in her score compared to the day before. (For example, Zelda's score on August 2 was $1,700 + 80 = 1,780$ points.) What was Zelda's average score in points over the $6$ days?
                    \begin{figure}[H]
                        \centering
                        \includegraphics[width=0.5\textwidth]{Problem/AMC_10_10A_2024_12_Problem_Figure_1.png}
                    \end{figure}
                    


 $\textbf{(A)}~1700\qquad\textbf{(B)}~1702\qquad\textbf{(C)}~1703\qquad\textbf{(D)}~1713\qquad\textbf{(E)}~1715$


 
 \textbf{Answer: }E



\textbf{Solution 1}

Going through the table, we see her scores over the six days were: $1700$, $1700+80=1780$, $1780-90=1690$, $1690-10=1680$, $1680+60=1740$, and $1740-40=1700$.

 Taking the average, we get  $\frac{(1700+1780+1690+1680+1740+1700)}{6} = \boxed{\textbf{(E) } 1715}.$

 -i\_am\_suk\_at\_math\_2

 



\textbf{Solution 2}

Compared to the first day $(1700)$, her scores change by $+80$, $-10$, $-20$, $+40$, and $+0$. So, the average is $1700 + \frac{80-10-20+40+0}{6} = \boxed{\textbf{(E) }1715}$.

 Sidenote: Many people would just add up the changes, which would give them 1700. That is wrong.

 -mathfun2012

 -JK2030 minor edit

 



\textbf{Solution 3}

As the scores of each day are dependent on previous days, we get: $1700 + \dfrac{0\cdot6 + 80\cdot5 + (-90)\cdot4 + (-10)\cdot3 + 60\cdot2 + (-40)\cdot1}{6} = \boxed{\textbf{(E) }1715}$

 



\textbf{Solution 4}

Plugging in x as 1700, we get the variables as: x + 80, x-10, x-20, x+40, and x.Finding the average of this gets us (6x + 90)/6 = x + 15Since x is 1700 we get our answer as 1700 + 15 = $\boxed{\textbf{(E) }1715}$

 


	\item (\textit{AMC 10 10A 2024}) Two transformations are said to commute if applying the first followed by the second
gives the same result as applying the second followed by the first. Consider these
four transformations of the coordinate plane:


 
                \begin{itemize}
                    \item a translation $2$ units to the right,
                \end{itemize}


                
                \begin{itemize}
                    \item a $90^{\circ}$-rotation counterclockwise about the origin,
                \end{itemize}


                
                \begin{itemize}
                    \item a reflection across the $x$-axis, and
                \end{itemize}


                
                \begin{itemize}
                    \item a dilation centered at the origin with scale factor $2.$
                \end{itemize}


                Of the $6$ pairs of distinct transformations from this list, how many commute?


 $\textbf{(A)}~1\qquad\textbf{(B)}~2\qquad\textbf{(C)}~3\qquad\textbf{(D)}~4\qquad\textbf{(E)}~5$


 
 \textbf{Answer: }C



\textbf{Solution 1 (Generalized)}

Label the given transformations $T_1, T_2, T_3,$ and $T_4,$ respectively. The rules of transformations are:

 
                \begin{itemize}
                    \item $T_1: \ (x,y)\to(x+2,y)$
                \end{itemize}


                
                \begin{itemize}
                    \item $T_2: \ (x,y)\to(-y,x)$
                \end{itemize}


                
                \begin{itemize}
                    \item $T_3: \ (x,y)\to(x,-y)$
                \end{itemize}


                
                \begin{itemize}
                    \item $T_4: \ (x,y)\to(2x,2y)$
                \end{itemize}


                Note that:

 <ul style="list-style-type:square;"> <li>Applying $T_1$ and then $T_2$ gives $(x,y)\to(x+2,y)\to(-y,x+2).$ Applying $T_2$ and then $T_1$ gives $(x,y)\to(-y,x)\to(-y+2,x).$ 

 Therefore, $T_1$ and $T_2$ do not commute. One counterexample is the preimage $(0,0).$

 </li>

 \item Applying $T_1$ and then $T_3$ gives $(x,y)\to(x+2,y)\to(x+2,-y).$ Applying $T_3$ and then $T_1$ gives $(x,y)\to(x,-y)\to(x+2,-y).$ 

 Therefore, $T_1$ and $T_3$ commute. They form a \textbf{glide reflection}.

 

 <li>Applying $T_1$ and then $T_4$ gives $(x,y)\to(x+2,y)\to(2x+4,2y).$ Applying $T_4$ and then $T_1$ gives $(x,y)\to(2x,2y)\to(2x+2,2y).$ 

 Therefore, $T_1$ and $T_4$ do not commute. One counterexample is the preimage $(0,0).$

 </li>

 <li>Applying $T_2$ and then $T_3$ gives $(x,y)\to(-y,x)\to(-y,-x).$ Applying $T_3$ and then $T_2$ gives $(x,y)\to(x,-y)\to(y,x).$ 

 Therefore, $T_2$ and $T_3$ do not commute. One counterexample is the preimage $(1,0).$

 </li>

 \item Applying $T_2$ and then $T_4$ gives $(x,y)\to(-y,x)\to(-2y,2x).$ Applying $T_4$ and then $T_2$ gives $(x,y)\to(2x,2y)\to(-2y,2x).$ 

 Therefore, $T_2$ and $T_4$ commute.

 

 \item Applying $T_3$ and then $T_4$ gives $(x,y)\to(x,-y)\to(2x,-2y).$ Applying $T_4$ and then $T_3$ gives $(x,y)\to(2x,2y)\to(2x,-2y).$ 

 Therefore, $T_3$ and $T_4$ commute.

 

 </ul>Together, $\boxed{\textbf{(C)}~3}$ pairs of transformations commute: $T_1$ and $T_3, T_2$ and $T_4,$ and $T_3$ and $T_4.$

 <u>\textbf{Remark}</u>

 To show that two transformations do not commute, we only need one counterexample.

 



\textbf{Solution 2 (Specific)}

Label the transformations as follows:

 • a translation $2$ units to the right $(W)$

 • a $90^{\circ}$-rotation counterclockwise about the origin $(X)$

 • a reflection across the $x$-axis $(Y)$

 • a dilation centered at the origin with scale factor $2$ $(Z)$

 Now, examine each possible pair of transformations with the point $(1,0)$:

 $1.$ $W$ and $X$. $W\rightarrow X$ ends with the point $(0,3)$. Going $X\rightarrow W$ ends in the point $(1,2)$, so this pair does not work

 $2.$ $W$ and $Y$. $W\rightarrow Y$ gives the point $(3,0)$, and going $Y\rightarrow W$ ends in the same point. This pair is valid.

 $3.$ $W$ and $Z$. $W\rightarrow Z$ ends in the point $(6,0)$, while going the other way gives $(4,0)$. This pair isn't commute.

 $4.$ $X$ and $Y$. $X\rightarrow Y$. gives the point $(0,-1)$, while the other way gives $(0,1)$. Not a valid pair

 $5.$ $X$ and $Z$. $X\rightarrow Z$ ends in the point $(0,2)$, and $Z\rightarrow X$ also ends in $(0,2)$. This pair works.

 $6.$ $Y$ and $Z$. $Y\rightarrow Z$ gives the point $(2,0)$, and going the other way also ends in $(2,0)$. This pair is valid.

 Therefore, the answer is $\boxed{\textbf{(C) }3}$.

 Note: It is easier to just visualize this problem instead of actually calculating points on paper.

 



\textbf{Solution 3 (Specific)}

Label the transformations as follows:

 • a translation $2$ units to the right $(A)$

 • a $90^{\circ}$-rotation counterclockwise about the origin $(B)$

 • a reflection across the $x$-axis $(C)$

 • a dilation centered at the origin with scale factor $2$ $(D)$

 Now, we count each transformation individually. It is not hard to see that $AC, BD,$ and $CD$ are commutative (an easy way to test commutativity for some cases would be to have the original point on the $x$-axis). 

 In total, $\boxed{\textbf{(C) }3}$ transformation pairs commute.

 



\textbf{Solution 4 (Quick With Reasoning)}

Alright so realize that obviously rotation and dilation are commute, if you ever looked at the polar plane. Also, reflecting around the x-axis means that moving to the right two units and dilating are commute. You can easily see this by taking an arbitrary point and applying the transformations. The rest don’t work since you can easily test out using random points on the plane, which means the answer is $\boxed{\textbf{(C) }3}$.

 


	\item (\textit{AMC 10 10A 2024}) One side of an equilateral triangle of height $24$ lies on line $\ell$. A circle of radius $12$ is tangent to line $\ell$ and is externally tangent to the triangle. The area of the region exterior to the triangle and the circle and bounded by the triangle, the circle, and line $\ell$ can be written as $a \sqrt{b} - c \pi$, where $a$, $b$, and $c$ are positive integers and $b$ is not divisible by the square of any prime. What is $a + b + c$?  


 $\textbf{(A)}~72\qquad\textbf{(B)}~73\qquad\textbf{(C)}~74\qquad\textbf{(D)}~75\qquad\textbf{(E)}~76$


 
 \textbf{Answer: }D



\textbf{Solution 1}

Call the bottom vertices of the triangle $B$ and $C$ (the one closer to the circle is $C$) and the top vertex $A$. The tangency point between the circle and the side of the triangle is $D$, and the tangency point on line $\ell$ $E$, and the center of the circle is $O$.

 
Draw radii to the tangency points, the arc is $60$ degrees because $\angle ACB$ is $60$, and since $\angle DCE$ is supplementary, it's $120^{\circ}$.(Using Angle of Intersecting Secants Theorem) The sum of the angles in a quadrilateral is $360$, which means $\angle DOE$ is $60^{\circ}$

 
Triangle $ODC$ is $30$-$60$-$90$ triangle so CD is $4\sqrt{3}$. Since we have $2$ congruent triangles ($\triangle ODC$ and $\triangle OEC$), the combined area of both is $48\sqrt{3}$.  The area of the arc is $144 \cdot \frac{60}{360} \cdot \pi$ which is $24\pi$, so the area of the region is 
 \[48\sqrt{3}-24\pi\].

 

 \[a+b+c \text{ is } 48+3+24 \text{, which is } \boxed{\mathbf{(D)}~75}\]

 edits by KR

 \textbf{Note} There were two possible configurations from this problem; the one described in the solution above and the configuration in which the circle is tangent to the bottom of line $\ell$ and the base of the equilateral triangle. However, since the area in this configuration is simply $0,$ we can infer that the problem is talking about the configuration in Solution 1.

 



\textbf{Solution 2 (Strategic Triangles)}


                    \begin{figure}[H]
                        \centering
                        \includegraphics[width=0.5\textwidth]{Solution/AMC_10_10A_2024_14_Solution_Figure_1.png}
                    \end{figure}
                    For context, we want to find the area of the quadrilateral $BDCE$ and subtract the area of circular sector $BDE$.

 We are given that triangle $AMC$ is equilateral with height 24.

 Through $30-60-90$ triangles, we see that length
 \[MF = \frac{24}{\sqrt{3}} = 8\sqrt{3}\]Therefore the side length of $AMC$ is simply $16\sqrt{3}$.

 We now pay attention to the circle. Notice how because $AMC$ is equilateral, $\angle DCE$ is $120^\circ$.

 Notice how arc $DE$ is inscribed in both $\angle DCE$ and $\angle DBE$, therefore angles $\angle DCE$ and $\angle DBE$ are supplementary, and angle $DBE = 180^\circ - 120^\circ = 60^\circ$.

 We have that $BD = BE = 12$. we see that $\angle BDE = \angle BED = x^\circ$. We have that $\angle DBE = 60^\circ$. Therefore,
 \[2x + 60^\circ = 180^\circ\]
 \[2x=120^\circ\]
 \[x=60^\circ\]Therefore, triangle $BDE$ is a equilateral triangle with side length 12.

 Additionally, the area of sector BDE is $\frac{1}{6} \cdot (12)^2 \cdot \pi$, or $24\pi$.

 Because $BD=BE$, $\angle DCE + \angle DBE$, $\angle BDC = \angle BEC = 90^\circ$ (radius-chord theorem), and $\angle BDC + \angle BEC = 180^\circ$, quadrilateral $BDCE$ is a kite. The area of a kite is simply the diagonals $d_1$ and $d_2$ multiplied together and divided by 2.

 In this case our diagonals are DE and CB. We already know DE = 12, but to find CB, we construct a triangle $ABC$, with $AC = 16\sqrt{3}$The diagonals of a kite bisect their respective angles. Therefore, $\angle ACB = \frac{120^\circ}{2} = 60^\circ$.

 Then, triangles $ABC$ and $BDC$ are similar through a common angle and common side (with the angle of course in-between the side).

 Applying the same logic to triangle $AFC$, one can see that quadrilateral AFCB is also a kite that is similar to kite BDCE, and therefore $\angle ABC = 90^\circ$.

 Through $30-60-90$ triangles, $BC = \frac{1}{2}AC$, or $8\sqrt{3}$.

 The area of kite BCDE is $[BCDE] = \frac{8\sqrt{3} \cdot 12}{2}$, or $48\sqrt{3}$.

 We need $[BCDE] - [BDE]$, or $48\sqrt{3} - 24\pi$.

 This is in the form $a\sqrt{b} - c\pi$, with $a=48$, $b=3$, and $c=24$.

 We need $a+b+c$, which is $48+3+24 =$ $\boxed{\textbf{(D)}75}$

 



\textbf{Solution 3 (Quick Guess)}

This is kind of a 

 Since this problem involves equilateral triangles, the only possible number under the square root is $3$. Now subtracting all of the answer choices by $3$, we get:
 \[\textbf{(A)}~72-3=69\qquad\textbf{(B)}~73-3=70\qquad\textbf{(C)}~74-3=71\qquad\textbf{(D)}~75-3=72\qquad\textbf{(E)}~76-3=73\]

 Due to the even parity of the problem, we can safely assume that the answer is either $B$ or $D$, but as $D$ is a multiple of $12$ and $24$, we get the answer of $\boxed{\textbf{(D)}~75}$.

 



\textbf{Solution 4}

(pardon the diagrams :D)

 say the area we want to find is x.

 since the equilateral triangle has an internal angle of 60, the exterior angle formed by the triangle and the line is 120. simplifying the diagram you will get:

 \begin{verbatim}
\   \#\#\#\#\# \ \#\#\#\#\#\#\#\#  \#\#\#\#\#\#\#\#   \#\#\#\#\#\_\_\_\_\_\_\_\_\_\_\_
\end{verbatim}

make three of these so that each circle is tangent to the other 2 circles

 \begin{verbatim}
\    \    \#\#\#\#\#     \ \#\#\#\#\#\#\#\# \#\#\#\#\#\#\#\#\#\#\#\#\#\#\#\#\#\#\#\#\_\#\#\#\#\_\_\_\_\_\_\_\_\_\_\#\#\#\#\#\#\#/ \#\#\#\#\# \#\#\#\#\#/\#\#\#\#\#\#\#\#     / \#\#\#\#\#\#\#\#    /   \#\#\#\#\#\#
\end{verbatim}

Since they are 3 congruent triangles, you can make an equilateral triangle using their radius(12), with each vertex at the center of each circle. This will make an equilateral triangle of side length 24. if you look now, the area within the equilateral triangle consists of 3 $60/360$ of a circle, and 3 of x.

 we first find the area of the triangle, which is $24 * 12\sqrt{3}$, we then find the area of $60/360$ of a circle, which is $60/360 * 12^2\pi$, we subtract $24 * 12\sqrt{3}$ by $60/360 * 12^2\pi$, and divide by 3, yielding the area of x.

 \begin{verbatim}
$24 * 12\sqrt{3}$ - $60/360 * 12^2\pi$
\end{verbatim}

\_\_\_\_\_\_\_\_\_\_\_\_\_\_\_\_\_\_\_\_\_\_\_\_\_\_\_\_\_\_\_\_\_\_\_\_\_\_\_\_\_ = $48\sqrt{3}-24\pi$

 \begin{verbatim}
				3
\end{verbatim}


$a+b+c$ is $48+3+24$ which is $\boxed{\textbf{(D)}~75}$

 



\textbf{Solution 5 (oh no)}

Setting up the problem graphically:

 Define the parametric function of the circle \(r(\theta)\), centered at \((0,12)\) with radius \(12\), as:
 \[r(\theta) \;=\; \bigl(12\cos(\theta),\,12\sin(\theta) + 12\bigr).\] 

 The line \(\ell\) is taken to be the \(x\)-axis, so \(y = 0\). The circle is tangent to \(\ell\) at \((0,0)\).

 We only care about the side of the triangle that is also tangent to the circle; call this line \(f(x)\). Since the triangle is equilateral of height 24, one vertex is at \((0,24)\) (directly above the tangent point), but for now we focus on the particular side tangent to the circle.

 Because we expect that line to make a \(330^\circ\) angle with the positive \(x\)-axis (or a slope of \(-\sqrt{3}\)), let 
 \[f(x) \;=\; -\sqrt{3}\,x \;+\; c.\]

 We also know the perpendicular distance from the circle’s center \((0,12)\) to \(f(x)\) must be \(12\), because the circle of radius \(12\) is externally tangent to that line.  The distance from a point \(\bigl(x_0,y_0\bigr)\) to a line \(Ax + By + C = 0\) is: 

 
 \[\frac{|Ax_0 + By_0 + C|}{\sqrt{A^2 + B^2}}.\]

 Here, \(f(x) = -\sqrt{3}\,x + c\) can be rewritten as \(\sqrt{3}\,x + y - c = 0.\)  Plugging in \((x_0,y_0)=(0,12)\) and setting the distance to 12 yields:

 
 \[\frac{|\sqrt{3}\cdot 0 \;+\; 12 \;-\; c|}{\sqrt{(\sqrt{3})^2 + 1^2}}\;=\;12.\]

 Solving this and noting that \(f(x)\) has a negative \(y\)-intercept gives:

 
 \[f(x) \;=\; -\sqrt{3}\,x \;-\; 12.\]

 Next, we find the point at which \(r(\theta)\) intersects \(f(x)\).  Substituting \(r(\theta) \;=\; (12\cos(\theta), \,12\sin(\theta) + 12)\) into \(f(x) = -\sqrt{3}\,x - 12,\) we set:

 
 \[12\sin(\theta) + 12 \;=\; -\sqrt{3}\,\bigl(12\cos(\theta)\bigr) - 12.\]

 One helpful approach is to note that a line perpendicular to \(f(x)\) and passing through \((0,12)\) has slope \(1/\sqrt{3}\), so that line is \(g(x) = \frac{x}{\sqrt{3}} + 12.\)  Intersecting \(g(x)\) with \(f(x)\) gives:

 
 \[\frac{x}{\sqrt{3}} + 12 \;=\; -\sqrt{3}\,x - 12,\]

 which solves to \(x = -6\sqrt{3}\).  Then \(y = 6.\)  So the intersection point of the circle and the line is \(\bigl(-6\sqrt{3},\,6\bigr)\).

 To find the corresponding \(\theta\), use: 
 \[12\cos(\theta) = -6\sqrt{3}, \quad12\sin(\theta) + 12 = 6.\]

 From the first, \(\cos(\theta) = -\sqrt{3}/2\).  This corresponds to \(\theta = 5\pi/6\) or \(7\pi/6.\)  Checking the \(y\)-coordinate shows \(\theta = 7\pi/6\) indeed yields \(\sin(7\pi/6) = -1/2\), so \(12\sin(7\pi/6) + 12 = 6.\)

 The circle also meets \(\ell\) (the \(x\)-axis) at \((0,0)\).  In polar terms, that is \(\theta = 3\pi/2\).  We want the area under the circle’s parametric curve from \(\theta=7\pi/6\) to \(\theta=3\pi/2.\)

 Parametrically, the area under \(y(\theta)\) from \(\theta_1\) to \(\theta_2\) can be found via:

 
 \[\int_{\theta_1}^{\theta_2} y(\theta)\,x'(\theta)\,d\theta.\]

 Here \(x(\theta) = 12\cos(\theta)\), so \(x'(\theta) = -12\sin(\theta)\).  Hence the area under the circle from \(\theta = 7\pi/6\) to \(\theta = 3\pi/2\) is:

 
 \[\int_{7\pi/6}^{3\pi/2} \bigl(12\sin(\theta) + 12\bigr) \,\bigl(-12\sin(\theta)\bigr)\,d\theta.\]

 This simplifies (pulling out constants) to:

 
 \[-12 \int_{7\pi/6}^{3\pi/2}\bigl(\sin^2(\theta) + \sin(\theta)\bigr)\, d\theta\;=\;-12 \left[\int_{7\pi/6}^{3\pi/2} \sin^2(\theta)\,d\theta\;+\;\int_{7\pi/6}^{3\pi/2} \sin(\theta)\,d\theta\right].\]

 Evaluating these integrals yields:

 
 \[54\sqrt{3} \;-\; 24\pi.\]

 Next, find where \(f(x)=0\):

 
 \[-\sqrt{3}\,x \;-\; 12 = 0\quad\Longrightarrow\quad x = -\frac{12}{\sqrt{3}} \;=\; -4\sqrt{3}.\]

 However, we also need the relevant intersection interval for the region.  The line and the circle intersect at \(x=-6\sqrt{3}\).  So the area under \(f(x)\) from \(x=-6\sqrt{3}\) to \(x=-4\sqrt{3}\) is:

 
 \[\int_{x=-6\sqrt{3}}^{-4\sqrt{3}} \bigl(-\sqrt{3}\,x - 12\bigr)\,dx.\]

 Carrying out that integral (or carefully checking the geometry) gives \(6\sqrt{3}.\)

 Subtracting this triangular “cap” (\(6\sqrt{3}\)) from the circle’s sector area (\(54\sqrt{3}-24\pi\)) gives the final region of interest:

 
 \[\bigl(54\sqrt{3} - 24\pi\bigr) - 6\sqrt{3}\;=\;48\sqrt{3} - 24\pi.\]

 Hence the area can be written as

 
 \[a\sqrt{b} \;-\; c\pi\;=\;48\sqrt{3} - 24\pi,\]

 so \(a=48\), \(b=3\), and \(c=24.\)  Therefore,
 \[a + b + c \;=\; 48 + 3 + 24 \;=\; 75.\]

 \(\boxed{D)75}\)

 P.S.: Please don’t whip out calculus like this on an AMC 10.Yes, it’s “doable” without a calculator (I speak from painful experience).Yes, I did it, but it took forever, and trust me, nobody at the AMC office expects calculus.Also, this is a terriBLE method

 



\textbf{Solution 6 (Pure Geometry Approach — Region DCE minus Circular Segment)}


                \begin{itemize}
                    \item Step 1: Area of triangle $DCE$
                \end{itemize}


                - Given: $DE = 12$ (radius of the circle), and the corresponding height from $C$ to $DE$ is $2\sqrt{3}$.- The area formula for a triangle is:

 \begin{verbatim}

 \[S_{\triangle DCE} = \frac{1}{2} \times \text{base} \times \text{height} = \frac{1}{2} \times 12 \times 2\sqrt{3} = 12\sqrt{3}\]
\end{verbatim}

- Explanation: This step finds the area of triangle $DCE$ using its base and height.

 ---

 
                \begin{itemize}
                    \item Step 2: Area of sector $ODE$
                \end{itemize}


                - The radius is $12$, and the central angle is $60^\circ$.- The formula for the area of a sector is:

 \begin{verbatim}

 \[S_{\text{sector ODE}} = \frac{60}{360} \times \pi \times 12^2 = 24\pi\]
\end{verbatim}

- Explanation: The sector's area is a fraction ($\frac{1}{6}$) of the full circle's area.

 ---

 
                \begin{itemize}
                    \item Step 3: Area of triangle $ODE$
                \end{itemize}


                
- $OD = OE = 12$, and the included angle $\angle DOE = 60^\circ$.- The formula for an equilateral triangle's area:

 \begin{verbatim}

 \[S_{\triangle ODE} = \frac{\sqrt{3}}{4} \times 12^2 = 36\sqrt{3}\]
\end{verbatim}

- Explanation: Triangle $ODE$ is equilateral, so we use the standard area formula for such a triangle.

 ---

 
                \begin{itemize}
                    \item Step 4: Area of the circular segment (arc $DE$ minus chord $DE$)
                \end{itemize}


                \begin{verbatim}

 \[S_{\text{segment}} = S_{\text{sector}} - S_{\triangle ODE} = 24\pi - 36\sqrt{3}\]
\end{verbatim}

- Explanation: The segment is the region between the arc and the chord $DE$.

 ---

 
                \begin{itemize}
                    \item Step 5: Final required region
                \end{itemize}


                \begin{verbatim}

 \[S_{\text{required}} = S_{\triangle DCE} - S_{\text{segment}}   = 12\sqrt{3} - (24\pi - 36\sqrt{3})   = 12\sqrt{3} - 24\pi + 36\sqrt{3}   = 48\sqrt{3} - 24\pi\]
\end{verbatim}

- Explanation: The required area is the triangle $DCE$ minus the circular segment. This gives the area in the desired region bounded by the triangle, the circle, and the line.

 ---

 
                \begin{itemize}
                    \item Step 6: Answer in the form $a\sqrt{b} - c\pi$
                \end{itemize}


                \begin{verbatim}

 \[a = 48, \quad b = 3, \quad c = 24\] 
 \[a + b + c = \boxed{75}\]
\end{verbatim}


                \begin{itemize}
                    \item Final Answer:
                \end{itemize}


                $48\sqrt{3} - 24\pi$, and $a + b + c = \boxed{75}$

 -clicksong

 



\textbf{Solution 7}


                \begin{itemize}
                    \item Step 1: Area of triangle $DCE$ (using Heron's formula)*
                \end{itemize}


                - By geometry, $\triangle DCE$ is a triangle with sides $DC = 12$, $EC = 12$, $DE = 12\sqrt{3}$.- Let $a = 12$, $b = 12$, $c = 12\sqrt{3}$.- The semi-perimeter is:

 \begin{verbatim}

 \[s = \frac{a + b + c}{2} = \frac{12 + 12 + 12\sqrt{3}}{2} = 12 + 6\sqrt{3}\]
\end{verbatim}

- The area by Heron's formula:

 \begin{verbatim}

 \[S_{\triangle DCE} = \sqrt{s(s-a)(s-b)(s-c)}\] 
 \[= \sqrt{(12+6\sqrt{3})(6\sqrt{3})(6\sqrt{3})(12-6\sqrt{3})}\] 
 \[= \sqrt{(12+6\sqrt{3})(12-6\sqrt{3})(6\sqrt{3})^2}\] 
 \[= \sqrt{(144 - 108) \times 108}   = \sqrt{36 \times 108}   = \sqrt{3888}   = 36\sqrt{3}\]
\end{verbatim}

- However, the actual required region is only $\frac{4}{3}$ of this triangle (since the shaded region is less than the whole triangle).- So,

 \begin{verbatim}

 \[S_{\triangle DCE,\text{needed}} = \frac{36\sqrt{3}}{1.5} = 24\sqrt{3}\]
\end{verbatim}

- But based on precise geometry (as in AMC standard), the correct area to use for this region is $12\sqrt{3}$ (check the triangle's effective area).

 ---

 
                \begin{itemize}
                    \item Step 2: Area of sector $ODE$*
                \end{itemize}


                - The radius is $12$ and the central angle is $60^\circ$.- The formula for the area of a sector is:

 \begin{verbatim}

 \[S_{\text{sector ODE}} = \frac{60}{360} \times \pi \times 12^2 = 24\pi\]
\end{verbatim}

---

 
                \begin{itemize}
                    \item Step 3: Area of triangle $ODE$*
                \end{itemize}


                - $OD = OE = 12$, $\angle DOE = 60^\circ$.- The formula for an equilateral triangle's area:

 \begin{verbatim}

 \[S_{\triangle ODE} = \frac{\sqrt{3}}{4} \times 12^2 = 36\sqrt{3}\]
\end{verbatim}

---

 
                \begin{itemize}
                    \item Step 4: Area of the circular segment (arc $DE$ minus chord $DE$)*
                \end{itemize}


                \begin{verbatim}

 \[S_{\text{segment}} = S_{\text{sector}} - S_{\triangle ODE} = 24\pi - 36\sqrt{3}\]
\end{verbatim}

---

 
                \begin{itemize}
                    \item Step 5: Final required region*
                \end{itemize}


                \begin{verbatim}

 \[S_{\text{required}} = S_{\triangle DCE} - S_{\text{segment}}   = 12\sqrt{3} - (24\pi - 36\sqrt{3})   = 12\sqrt{3} - 24\pi + 36\sqrt{3}   = 48\sqrt{3} - 24\pi\]
\end{verbatim}

- Explanation: The required area is the triangle $DCE$ minus the circular segment.

 ---

 
                \begin{itemize}
                    \item Step 6: Expressing the answer as $a\sqrt{b} - c\pi$*
                \end{itemize}


                \begin{verbatim}

 \[a = 48, \quad b = 3, \quad c = 24\] 
 \[a + b + c = \boxed{75}\]
\end{verbatim}


                \begin{itemize}
                    \item Final Answer:*
                \end{itemize}


                $48\sqrt{3} - 24\pi$, and $a + b + c = \boxed{75}$

 also by clicksong

 



\textbf{Solution 8 - HEFLE (High Effort For Low Efficiency) a.k.a. Calculus Raigebait}


                    \begin{figure}[H]
                        \centering
                        \includegraphics[width=0.5\textwidth]{Solution/AMC_10_10A_2024_14_Solution_Figure_2.png}
                    \end{figure}
                    

 So, we're here to over-complicate things. Let's start with the circle. Recognize that since the triangle is equilateral, all angles measure $60^\circ$ Observe that the two tangent lines form a $120^\circ$ angle due to the Angle Sum Theorem. Therefore, the area within the circle is the area within a $60^\circ$ arc.

 Let's complicate finding the area of the circle.Now, instead of using the equation of a circle we set 
 \[x^2 + y^2 = 144\]

 then we solve for $y$: 
 \[y = \pm\sqrt{144 - x^2}\] then we want the bottom of this circle so we do the negative root $y$ = $-\sqrt{144 - x^2}$ Now, we define this as a function $f(x)$

 
 \[f(x) = -\sqrt{144 - x^2}\] but hold up! waaait a minute! We can make the circle centered at $(12,0)$ because we're just like that. So, to translate it 12 right we must do $f(x - 12)$ so 
 \[f(x - 12) = -\sqrt{144 - (x - 12)^2} = -\sqrt{-x^2 + 24x}\] and because there is a $60 ^\circ$  arc we want $\frac{1}{3}$ of this semicircle so we compute 
 \[\frac{1}{3} \cdot \int_{0}^{24} -\sqrt{-x^2 + 24x} \ \  dx\] then we need to undo our steps and negate the integral(this is done when it's converted to polar) 
 \[\frac{1}{3} \cdot \int_{0}^{24} -\sqrt{-x^2 + 24x} \ \ dx \rightarrow \frac{1}{3} \cdot \int_{0}^{24} -\sqrt{144-(x-12)^2}\ \ dx \rightarrow \int_0^{\pi/3} \frac{1}{2} \cdot 12^2  \ d\theta = \int_0^{\pi/3} 72 \ d\theta\]Recall that
 \[A= \frac{1}{2} \cdot \int_{\theta_1}^{\theta_2} [f(\theta)]^2 \  \ d\theta\]Therefore,
 \[\int_0^{\pi/3} 72 \  d\theta = 72 \cdot \frac{\pi}{3} = 24\pi\]Now, we must find the area enclosed by $2$ radii extending to the two points of tangency and their respective tangent lines. So, the two radii are $12$, and they form a  $60^\circ$ angle. Therefore, you have an equilateral triangle of side length $12$, and its height is $6\sqrt{3}$ using $30^\circ \ 60^\circ \ 90^\circ$ 

 Now, you can just see that the area is $36\sqrt{3}$, but we won't be doing that here. We can define an absolute value function $g(x)$ such that $g(x)$ has zeros at $-6$ and $6$, and a y-intercept at $(0, 6\sqrt{3})$. This would create our triangle above the $x$-axis. We can see that its slope is $\sqrt{3}$ for $x<0$ and $-\sqrt{3}$ for $x>0$. Then, we can see that it's translated $6\sqrt{3}$ up from its parent function $h(x)=|x|$. Therefore, the function is $g(x) = -|x\sqrt{3}|+6\sqrt{3}$. Now, we integrate over $x\in[-6, 6]$
 \[\int_{-6}^{6} -|x\sqrt{3}|+6\sqrt{3} \ \  dx\]Now, we split this into a piecewise function over two integrals
 \[\int_{-6}^{0} x\sqrt{3}+6\sqrt{3} \ \  dx \ + \ \int_{0}^{6} -x\sqrt{3}+6\sqrt{3} \ \  dx\]Let's apply the power rule to each term. We get
 \[\int_{-6}^{0} x\sqrt{3}+6\sqrt{3} \ \ dx= \left[\frac{x^2\sqrt3}{2}+6x\sqrt{3} \right]_{-6}^0 = \left(\frac{0^2\sqrt3}{2}+6\cdot0\sqrt{3}\right) - \left(\frac{(-6)^2\sqrt3}{2}+6\cdot-6\sqrt{3}\right) =\]
 \[0 - (18\sqrt{3} - 36\sqrt3) = 0 - (-18\sqrt{3}) = 18\sqrt{3}\]and, we could say that they are mirror images of each other, but that would make this faster (it's called inefficient for a reason)
 \[\int_{0}^{6} -x\sqrt{3}+6\sqrt{3} \ \  dx = \left[\frac{-x^2\sqrt3}{2}+6x\sqrt{3} \right]_{0}^6 = \left(\frac{-6^2\sqrt3}{2}+6\cdot6\sqrt{3}\right) - \left(\frac{(0)^2\sqrt3}{2}+6\cdot0\sqrt{3}\right) =\]
 \[(-18\sqrt{3}+36\sqrt{3}) - 0 = 18\sqrt{3}\]Therefore, by the substitution property equality,
 \[\int_{-6}^{0} x\sqrt{3}+6\sqrt{3} \ \  dx \ + \ \int_{0}^{6} -x\sqrt{3}+6\sqrt{3} \ \  dx = 18\sqrt{3}+18\sqrt{3}=36\sqrt{3}\]Then, there is another triangle we must define above the $x$-axis. It is enclosed by the equilateral triangle of height $24$, line $\ell$, and the triangle we just found the area of. This triangle has base $12$ and height of $\frac{6}{\sqrt{3}}$ using $30^\circ \ 60^\circ \ 90^\circ$ triangle ratios, which rationalizes to $2\sqrt{3}$. So, you can tell the area is $12\sqrt{3}$, but that would make our solution faster.

 Since half the base is $6$, we can tell the zeros are $x= -6$ and $x=6$, and the $y$-intercept is $2\sqrt{3}$ 

 We can tell the slope is $\frac{\sqrt{3}}{3}$ for $x<0$ and $-\frac{\sqrt{3}}{3}$ for $x>0$

 Then, we define another function $f(x)= -|\frac{1}{3} \cdot x\sqrt{3}| + 2\sqrt{3}$

 Now, we integrate over $x\in[-6, 6]$So, we must compute 
 \[\int_{-6}^{6} -|\frac{1}{3} \cdot x\sqrt{3}| + 2\sqrt{3} \ dx\]We split the function into a piecewise over two integrals
 \[\int_{-6}^{0} \ \frac{1}{3} \cdot x\sqrt{3} + 2\sqrt{3} \ dx + \int_{0}^{6} -\frac{1}{3} \cdot x\sqrt{3} + 2\sqrt{3} \ \ dx\]Now, we apply the reverse power rule to each term, we get
 \[\int_{-6}^{0} \  \frac{1}{3} \cdot x\sqrt{3} + 2\sqrt{3} \ \ dx = \left[ \frac{\sqrt{3}}{{3}} \cdot \frac{x^2}{2}+2x\sqrt{3} \right]_{-6}^{0} =\left[ \frac{\sqrt{3}}{{6}} \cdot{x^2}+2x\sqrt{3} \right]_{-6}^{0} =\]
 \[\left(\frac{\sqrt{3}}{{6}} \cdot{0^2}+2\cdot0\sqrt{3}\right)-\left(\frac{\sqrt{3}}{{6}} \cdot{6^2}+2\cdot-6\sqrt{3}\right)= 0-\left(\frac{\sqrt{3}}{{6}} \cdot{6^2}+2\cdot-6\sqrt{3}\right) =\]
 \[0- (6\sqrt{3}-12\sqrt{3})=6\sqrt{3}\]Then
 \[\int_{0}^{6} -\frac{1}{3} \cdot x\sqrt{3} + 2\sqrt{3} \ \ dx = \left[ -\frac{\sqrt{3}}{{3}} \cdot \frac{x^2}{2}+2x\sqrt{3} \right]_{0}^{6} = \left[ -\frac{\sqrt{3}}{{6}} \cdot{x^2}+2x\sqrt{3} \right]_{0}^{6}  =\]
 \[\left(-\frac{\sqrt{3}}{{6}} \cdot{6^2}+2\cdot6\sqrt{3}\right) - \left(-\frac{\sqrt{3}}{{6}} \cdot{0^2}+2\cdot0\sqrt{3} \right)=(-6\sqrt{3} + 12\sqrt{3})-0= 6\sqrt{3}\]Now, by the substitution property of equality,
 \[\int_{-6}^{0} \frac{1}{3} \cdot x\sqrt{3} + 2\sqrt{3} \ \ dx + \int_{0}^{6} -\frac{1}{3} \cdot x\sqrt{3} + 2\sqrt{3} \ \ dx = 6\sqrt{3}+6\sqrt{3}=12\sqrt{3}\]Now, the total area enclosed by the two circle radii of length $12$, line $\ell$, and the isosceles triangle of height $24$ is simply
 \[36\sqrt{3}+12\sqrt{3}=48\sqrt{3}\]Now, we must subtract the circle's sector area of $24\pi$ to get the expression
 \[48\sqrt{3}-24\pi\]Which is in the form 
 \[a\sqrt{b}-c\pi\]That the problem asks for. Hence,
 \[a=48\]
 \[b=3\]
 \[c=24\]The problem asks for $a+b+c$, thus
 \[a+b+c=48+3+24=75 \implies \boxed{\textbf{(D)}~75}\]-shockfront99

 PS: I've never taken calculus, so I could be pretty wrong

 


	\item (\textit{AMC 10 10A 2024}) Let $M$ be the greatest integer such that both $M+1213$ and $M+3773$ are perfect squares. What is the units digit of $M$?


 $\textbf{(A) }1\qquad\textbf{(B) }2\qquad\textbf{(C) }3\qquad\textbf{(D) }6\qquad\textbf{(E) }8$


 
 \textbf{Answer: }E



\textbf{Solution 1}

Let $M+1213=P^2$ and $M+3773=Q^2$ for some positive integers $P$ and $Q.$ We subtract the first equation from the second, then apply the difference of squares: 
 \[(Q+P)(Q-P)=2560.\] Note that $Q+P$ and $Q-P$ have the same parity, and $Q+P>Q-P.$ 

 We wish to maximize both $P$ and $Q,$ so we maximize $Q+P$ and minimize $Q-P.$ It follows that
 \begin{align*} Q+P&=1280, \\ Q-P&=2, \end{align*}
from which $(P,Q)=(639,641).$

 Finally, we get $M=P^2-1213=Q^2-3773\equiv1-3\equiv8\pmod{10},$ so the units digit of $M$ is $\boxed{\textbf{(E) }8}.$

 



\textbf{Solution 2}

Ideally, we would like for the two squares to be as close as possible. The best case is that they are consecutive squares (no square numbers in between them); however, since $M+1213$ and $M+3773$ (and thus their square roots) have the same parity, they cannot be consecutive squares (which are always opposite parities). The best we can do is that $M+1213$ and $M+3773$ have one square in between them. 

 Let the square between $M+1213$ and $M+3773$ be $x^2$. So, we have $M+1213 = (x-1)^2$ and $M+3773 = (x+1)^2$. Subtracting the two, we have $(M+3773)-(M+1213) = (x+1)^2 - (x-1)^2$, which yields $2560 = 4x$, which leads to $x = 640$. Therefore, the two squares are $639^2$ and $641^2$, which both have units digit $1$. Since both $1213$ and $3773$ have units digit $3$, $M$ will have units digit $\boxed{\textbf{(E) }8}$.

 



\textbf{Solution 3}

Let $M+1213=N^2$ $\Rightarrow M+3773=(N+a)^2$

 It is obvious that $a\neq1$ by parity

 Thus, the minimum value of $a$ is 2Which gives us,
 \[(N+a)^2-N^2=M+3773-(M+1213)\]
 \[4N+4=2560\]
 \[N=639\]Plugging this back in,
 \[M=N^2-1213 \space \mod \space 10\]
 \[M=8 \space \mod \space 10\]Hence the answer $\boxed{\textbf{(E) }8}$.

 - trevian1(minor edit)

 



\textbf{Solution 4}

Let $M+1213=n^2$ and $M+3773=(n+1)^2$ for some positive integer $n$. We do this because, in order to maximize $M$, the perfect squares need to be as close to each other as possible. We find that this configuration doesn't work, as when we subtract the equations, we have $2n+1=2560$; impossible. Then we try $M+3773=(n+2)^2$. Now we would have $4n+4=2560$ which indeed works! $n=639$. 

 Finally, we get $M=n^2-1213$ so the units digit of $M$ is $11-3=\boxed{\textbf{(E) }8}.$

 



\textbf{Solution 5 (Quadratic Residues, Answer Choices rather than finding )}

Every perfect square modulo 4 must end in 0 or 1 modulo 4. Reducing modulo 4 yields 

 
 \[M + 1 \pmod{4}\]
 \[M + 1 \pmod{4}\].

 Both values $M$ yield the same result meaning that the units digit of $M$ is unchanged through constriction. Therefore, $M \equiv 3, 0 \pmod{4}$. Analyzing the answer choices, only 3 and 8 work.

 Now consider quadratic residue's modulo 7. Once again, they are either 0, 1, 2, 4 modulo 7. Reducing, we have

 
 \[M + 2 \pmod{7}\]
 \[M \pmod{7}\].

 If the units digit of $M$ is 3, we arise at a contradiction that $M+2 \pmod{7}$ can never satisfy a perfect square. Thus, $M \pmod{10} = \boxed{8}$.

 



\textbf{Solution 6 (Fakesolve)}

Squares can only end in 0,1,4,6,5, and 9. To have x+3 (since 1213 and 3773 end in 3) end in one of those numbers and be maximized, it would be 8, because 8+3=11. The answer is $\boxed{\textbf{(E) }8}.$

 
Note: This solution is just a lucky fake solve; please look at the other solutions. The note below this gives a good idea of why the solution is not meaningful. M is uniquely constrained by defined distances from two perfect squares. There is no "maximizing" M, you simply have to solve for it algebraically.

 \textbf{Note} We experiment with values of $M$ to find the reason for why $M$ is maximised when $M + 1213$ and $M + 3773$ are nearly consecutive perfect squares. If $M$ is very small, $M + 1213$ and $M + 3773$ are perfect squares that are far apart. Yet, as $M$ grows, the relative distance between $M + 1213$ and $M + 3773$ decreases, so for very nearly consecutive perfect squares, $M$ is very large.

 (I don't know if this is trivial - when I first read the solutions, I was confused to why this was true. Maybe I didn't get enough sleep.)

 


	\item (\textit{AMC 10 10A 2024}) All of the rectangles in the figure below, which is drawn to scale, are similar to the enclosing rectangle. Each number represents the area of the rectangle. What is length $AB$? $\newline$

                    \begin{figure}[H]
                        \centering
                        \includegraphics[width=0.5\textwidth]{Problem/AMC_10_10A_2024_16_Problem_Figure_1.png}
                    \end{figure}
                    
$\textbf{(A) }4+4\sqrt5\qquad\textbf{(B) }10\sqrt2\qquad\textbf{(C) }5+5\sqrt5\qquad\textbf{(D) }10\sqrt[4]{8}\qquad\textbf{(E) }20$


 
 \textbf{Answer: }D



\textbf{Solution 1}

Using the rectangle with area $1$, let its short side be $x$ and the long side be $y$. Observe that for every rectangle, since ratios of the side length of the rectangles are directly proportional to the ratios of the square roots of the areas (For example, each side of the rectangle with area $9$ is $\sqrt{9}=3$ times that of the rectangle with area $1$), as they are all similar to each other.

 The side opposite $AB$ on the large rectangle is hence written as $6x + 4x + 2y\sqrt{2} + 3y\sqrt{2} = 10x+5y\sqrt{2}$. However, $AB$ can be written as $4y\sqrt{2}+5x+7x = 4y\sqrt{2}+12x$. Since the two lengths are equal, we can write $10x+5y\sqrt{2} = 4y\sqrt{2}+12x$, or $y\sqrt{2} = 2x$. Therefore, we can write $y=x\sqrt{2}$.

 Since $xy=1$, we have $(x\sqrt{2})(x) = 1$, which we can evaluate $x$ as $x=\frac{1}{\sqrt[4]{2}}$. From this, we can plug back in to $xy=1$ to find $y=\sqrt[4]{2}$. Substituting into $AB$, we have $AB = 4y\sqrt{2}+12x = 4(\sqrt[4]{2})(\sqrt{2})+\frac{12}{\sqrt[4]{2}}$ which can be evaluated to $\boxed{\textbf{(D) }10\sqrt[4]{8}}$.

 \textbf{Remark} We know that the area is an integer, so after finding $y=x\sqrt{2}$, AB must contain a 4th root. The only such option is $\boxed{\textbf{(D) }10\sqrt[4]{8}}$.

 



\textbf{Solution 2}

Let the rectangle's height be $x,$ the length $AB=y.$ The entire rectangle has an area of $200.$ We will be using this fact for ratios.

 Note that the short side of the rectangle with area 32 will have a height of $\sqrt{\frac{32}{200}}\cdot x = \frac{2}{5}x.$ We use $x$ because it is apparent that the height of the rectangle with area $32$ is the shorter side, corresponding with $x.$

 Similarly, the long side of the rectangle with area 36 has a height of $\sqrt{\frac{36}{200}}\cdot y = \frac{3\sqrt{2}}{10}y.$

 Noting that the total height of the big rectangle has height $x,$ we have the equation $\frac25 x + \frac{3\sqrt{2}}{10}y = x \Rightarrow x=\frac{y}{\sqrt{2}}.$

 Since the area $xy=\frac{y^2}{\sqrt{2}}$ is equal to 200, we have:

 
 \begin{align*} y&=\sqrt{200\sqrt{2}} \\ &=\sqrt{100\sqrt{8}} \\ &=\boxed{\textbf{(D) }10\sqrt[4]{8}}. \end{align*}


 



\textbf{Solution 3 - ruler (last effort)}

Given that the diagram is drawn to scale, we can use a ruler to estimate the length of $AB$. 

 We start by measuring the lengths of the rectangle with area $1$, which may vary per viewing medium. For the sake of the solution, we use side lengths ~$\frac{7}{10}$ cm and ~$\frac{5}{10}$ cm. 

 To get the scale ratio from centimeters to the units in the problem, we need to find a ratio $x$ such that $\frac{7x}{10}\cdot\frac{5x}{10}=1$ 

 Solving this equation, we get $x=\frac{10}{\sqrt{35}}$

 We then measure the length of $AB$(will vary), to get ~$10.2$ cm. We multiply this length by our early ratio to get $\frac{10 \cdot 10.2}{\sqrt{35}} = \frac{102}{\sqrt{35}} \approx \frac{102}{6} \approx 17$

 The answer choice closest to this would be $10\sqrt[4]{8}\approx16.8$, so therefore the closest answer is $\boxed{\textbf{(D) }10\sqrt[4]{8}}$.

 Note:Never ever use this in an actual test unless you're giving it a last ditch attempt.

 \textbf{Remark} The specific numbers used in the solution above vary per test medium, but the method should still work.

 


	\item (\textit{AMC 10 10A 2024}) Two teams are in a best-two-out-of-three playoff: the teams will play at most $3$ games, and the winner of the playoff is the first team to win $2$ games. The first game is played on Team A's home field, and the remaining games are played on Team B's home field. Team A has a $\frac{2}{3}$ chance of winning at home, and its probability of winning when playing away from home is $p$. Outcomes of the games are independent. The probability that Team A wins the playoff is $\frac{1}{2}$. Then $p$ can be written in the form $\frac{1}{2}(m - \sqrt{n})$, where $m$ and $n$ are positive integers. What is $m+n$? 


 $\textbf{(A)}~10\qquad\textbf{(B)}~11\qquad\textbf{(C)}~12\qquad\textbf{(D)}~13\qquad\textbf{(E)}~14$


 
 \textbf{Answer: }E



\textbf{Solution}

We only have three cases where A wins: AA, ABA, and BAA (A denotes a team A win and B denotes a team B win). Knowing this, we can sum up the probability of each case. Thus the total probability is $\frac{2}{3}p+\frac{2}{3}(1-p)p+\frac{1}{3}p^2=\frac{1}{2}$. Multiplying  both sides by 6 yields $4p+4p(1-p)+2p^2=3$, so $2p^2-8p+3=0$ and we find that $p=\frac{4\pm\sqrt{10}}{2}$. Luckily, we know that the answer should contain  $\frac{1}{2}(m - \sqrt{n})$, so the solution is $p=\frac{4-\sqrt{10}}{2}=\frac{1}{2}(4-\sqrt{10})$ and the answer is $4+10=\boxed{\textbf{(E) } 14}$.

 Another way to see the answer is subtraction and not addition is to realize that $p$ is between $0$ and $1$ since it is a probability.~andliu766

 


	\item (\textit{AMC 10 10A 2024}) There are exactly $K$ positive integers $5 \leq b \leq 2024$ such that the base-$b$ integer $2024_{b}$ is divisible by $16$(where $16$ is in base ten). What is the sum of the digits of $K$?


 $\textbf{(A)}~16\qquad\textbf{(B)}~17\qquad\textbf{(C)}~18\qquad\textbf{(D)}~20\qquad\textbf{(E)}~21$


 
 \textbf{Answer: }D



\textbf{Solution 1}

$2024_b = 2b^3+2b+4\equiv 0\pmod{16}\implies b^3+b+2\equiv 0\pmod 8$. If $b$ is even, then $b+2\equiv 0\pmod 8\implies b\equiv 6\pmod 8$. If $b$ is odd, then $b^2\equiv 1\pmod 8$* $\implies b^3+b+2\equiv 2b+2\pmod 8$ so $2b+2\equiv 0\pmod 8\implies b+1\equiv 0\pmod 4\implies b\equiv 3,7\pmod 8$. Now $8\mid 2024$ so $\frac38\cdot 2024=759$ but one of the answers we got from that, $3$, is too small, so $759 - 1 = 758\implies\boxed{\textbf{(D) }20}$.

 
                \begin{itemize}
                    \item The reason why $b^2\equiv 1\pmod 8$ (if b is odd) is that we can write $b$ as $2k+1,$ where $k$ is any number (such that the problem conditions are met). Plugging this into $b^2\equiv 1\pmod 8,$ we get $(2k+1)^2=4k^2+1+4k=4k(k+1)+1\equiv 1\pmod 8.$ Since $4k$ is, of course, always a multiple of $4$ and since $k+1$ is even, $4k(k+1)$ is divisible by $8,$ and adding $1$ creates a remainder of $1.$
                \end{itemize}


                



\textbf{Solution 2}

\begin{align*}2024_b\equiv0\pmod{16} \\2b^3+2b+4\equiv0\pmod{16} \\b^3+b+2\equiv0\pmod8 \\\end{align*}


 Clearly, $b$ is either even or odd. If $b$ is even, let $b=2a$. 

 \begin{align*}(2a)^3+2a+2\equiv0\pmod8 \\8a^3+2a+2\equiv0\pmod8 \\0+2a+2\equiv0\pmod8 \\a+1\equiv0\pmod4 \\a\equiv3\pmod4 \\\end{align*}


 Thus, one solution is $b=2(4x+3)=8x+6$ for some integer $x$, or $b\equiv6\pmod8$. 

 What if $b$ is odd? Then let $b=2a+1$: 

 \begin{align*}(2a+1)^3+2a+1+2\equiv0\pmod8 \\8a^3+12a^2+6a+1+2a+1+2\equiv0\pmod8 \\8a^3+12a^2+8a+4\equiv0\pmod8 \\4a^2+4\equiv0\pmod8 \\a^2\equiv1\pmod2 \\\end{align*}


 This simply states that $a$ is odd. Thus, the other solution is $b=2(2x+1)+1=4x+3$ for some integer $x$, or $b\equiv3\pmod4$. 

 We now simply must count the number of integers between $5$ and $2024$, inclusive, that are $6$ mod $8$ or $3$ mod $4$. Note that the former case comprises even numbers only while the latter is only odd; thus, there is no overlap and we can safely count the number of each and add them. 

 In the former case, we have the numbers $6,14,22,30,\dots,2022$; this list is equivalent to $8,16,24,32,\dots,2024\cong1,2,3,4,\dots,253$, which comprises $253$ numbers. In the latter case, we have the numbers $7,11,15,19,\dots,2023\cong4,8,12,16,\dots,2020\cong1,2,3,4,\dots,505$, which comprises $505$ numbers. There are $758$ numbers in total, so our answer is $7+5+8=\boxed{\textbf{(D) 20}}$. 

 



\textbf{Solution 3}

Note that $2024_b=2b^3+2b+4$ is to be divisible by $16$, which means that $b^3+b+2$ is divisible by $8$.

 If $b=0$, then $b^3+b+2 \equiv (0)^3 + (0) + 2 \equiv 2$ is not divisible by $8$.

 If $b=1$, then $b^3+b+2 \equiv (1)^3 + (1) + 2  \equiv 4$ is not divisible by $8$.

 If $b=2$, then $b^3+b+2 \equiv (2)^3 + (2) + 2  \equiv 4$ is not divisible by $8$.

 If $b=3$, then $b^3+b+2 \equiv (3)^3 + (3) + 2  \equiv (8+1)\cdot3 + (3) + 2 \equiv 0$ is divisible by $8$.

 If $b=4$, then $b^3+b+2 \equiv (4)^3 + (4) + 2 \equiv 0 + 4 + 2 \equiv 6$ is not divisible by $8$.

 If $b=5$, then $b^3+b+2 \equiv (-3)^3 + (-3) + 2 \equiv  (8+1)\cdot 3 + (-3) + 2 \equiv 2$ is not divisible by $8$.

 If $b=6$, then $b^3+b+2 \equiv (-2)^3 + (-2) + 2 \equiv -8 + (-2) + 2 \equiv 0$ is divisible by $8$.

 If $b=7$, then $b^3+b+2 \equiv (-1)^3 + (-1) + 2 \equiv -1 + (-1) + 2 \equiv 0$ is divisible by $8$.

 Therefore, for every $8$ values of $b$, $3$ of them will make $b^3+b+2$ divisible by $8$. Therefore, since $2024$ is divisible by $8$, $\dfrac{3}{8}\cdot2024=759$ values of $b$, but this includes $b=3$, which does not satisfy the given inequality. Therefore, the answer is 
 \[759-1=758\rightarrow7+5+8=\boxed{\text{(D) }20}\] ~Tacos\_are\_yummy\_1

 More detail by ~\href{https://artofproblemsolving.com/wiki/index.php/User:Cyantist}{luckuso}

 



\textbf{Solution 4}

$2024_b=2\ast\ b^3+2\ast\ b+4\ \\ {2024}_{\left(b+8\right)}=2\ast\left(b+8\right)^3+2\ast\left(b+8\right)+4$${2024}_{\left(b+8\right)}-{2024}_b=2*\left(8\right)*\left(b^2+8b+64\right)+2*8\ =16*\left(b^2+8b+64\right)+16$

 \begin{align*}2024_{(b+8)}-2024_b\equiv0\ (mod\ 16)\\2024_{(b+8)}\ \ \equiv2024_b\ \ (mod\ 16)\\2024_0\equiv4\ (mod\ 16)\\2024_1\equiv8\ (mod\ 16)\\2024_2\equiv6\ (mod\ 16)\\2024_3\equiv0(mod\ 16)\\2024_4\equiv12(mod\ 16)\\2024_5\equiv8(mod\ 16)\\2024_6\equiv0(mod\ 16)\\2024_7\equiv0(mod\ 16)\\\end{align*}


 We need$b\ \equiv3\ (mod\ 8)\ \ or\ b\ \equiv6\ (mod\ 8)\ \ or\ b\ \equiv7\ (mod\ 8) \\ \lfloor(2024-3)/8\rfloor+\lfloor(2024-6)/8\rfloor+\lfloor(2024-7)/8\rfloor+3=759$take away one because $b=3$ is out of range, so $758\Rightarrow7+8+5=\boxed{\text{(D) }20}$

 



\textbf{Solution 4.1: Alternate Ending}

One can arrange 2024 into different "blocks" consisting of 8 numbers each, particularly \( \{0,1,2,3,4,5,6,7\} \), \( \{8,9,10,11,12,13,14,15\} \) ...

 These just represent different \( \mod(8) \) class residues. So, you can just divide 2024 by 8 blocks giving us 253 total blocks with 0 extra semi-blocks.

 We see that \( 2024_3 \mod 8 \), \( 2024_6 \mod 8 \), and \( 2024_7 \mod 8 \) all appear exactly 253 times each, so we have \( 253 \times 3 = 759 \) different values of \( K \).

 However, we must not include \( 2024_3 \), as the bases \( b \) range from \( 5 \le b \le 2024 \), therefore we have \( 759-1=758 \) cases.

 Our answer is \( 7+5+8= \) $\boxed{\text{(D) }20}$.

 



\textbf{Solution 5 (Factoring)}

Start similar to other solutions.

 
 \[b^3+b+2 \equiv 0\pmod 8\]

 Except now we notice the following:

 \begin{align*}b^3+8+b+2 \equiv 8\pmod 8 \equiv 0\pmod 8\\(b+2)(b^2-2b+4) + b+2 \equiv 0\pmod 8\\(b+2)(b^2-2b+5) \equiv 0\pmod 8\\\end{align*}


 Notice that $b+2$ and $b^2-2b+5$ have opposite parity, and when multiplied, is divisible by $8$. Thus, either $b+2$ or $b^2-2b+5$ is divisible by $8$.

 $b+2$ is obvious, where $b+2 | 8$ when $b \equiv 6\pmod 8$

 When testing for $b^2-2b+5 | 8$, we find $b \equiv 3, 7\pmod 8$

 Proceed similar to the other solutions, $\lfloor\dfrac{3}{8} \cdot 2024\rfloor-1=759-1=758$, where $7+5+8=\boxed{\text{(D) }20}$

 Also noting this is an example of a misplaced problem this would be far to easy for even P 18 of amc 10, deserved to be Q5-6

 



\textbf{Solution 6}

$2024_b =2b^3 +2b+4 \equiv 0 \pmod {16} \to 2b^3 +2b \equiv 12 \pmod {16}$. 

 Dividing everything by $2$ gives: 

 
 \[b^3+b \equiv 6 \pmod {\frac{16}{gcd(16,2)}} \to b(b^2 +1)\equiv 6 \pmod {8}\].

 Then proceed in the same manner as the solutions above to get a final answer of $\boxed{\textbf{(D) 20}}$. 

 


	\item (\textit{AMC 10 10A 2024}) The first three terms of a geometric sequence are the integers $a, 720$ and $b$, where $a < 720 < b$. What is the sum of the digits of the least possible value of $b$?


 $\textbf{(A) } 9\qquad\textbf{(B) } 12\qquad\textbf{(C) } 16\qquad\textbf{(D) } 18\qquad\textbf{(E) } 21$


 
 \textbf{Answer: }E



\textbf{Solution 1}

For a geometric sequence, we have $ab=720^2=2^8 3^4 5^2$, and we can test values for $b$. We find that $b=768$ and $a=675$ works, and we can test multiples of $5$ in between the two values. Finding that none of the multiples of 5 divide $720^2$ besides $720$ itself, we know that the answer is $7+6+8=\boxed{\textbf{(E) } 21}$.

 (Note: To find the value of $b$ without bashing, we can observe that $2^8=256$, and that multiplying it by $3$ gives us $768$, which is really close to $720$. ~ YTH)

 Note: The reason why $ab=720^2$ is because $b/720 = 720/a$. Rearranging this gives $ab = 720^2$

 Note: Another reason that $ab=720^2$ is because the $\sqrt{ab}=720$ (as the middle term in a geometric series is always the geometric mean [the geometric mean is the square root of the product of the first and last terms of the series]) and squaring on both sides results in $ab=720^2$.

 



\textbf{Solution 2}

$720 = 2^4 \cdot 3^2 \cdot 5$. We want to find factors $x$ and $y$ where $y>x$ such that $\frac{y}{x}$, the ratio, is minimized. To do this, $x$ should be as large as possible, and $y$ should be just slightly larger than $x$. Additionally, to ensure $a$ and $b$ are both integers, $x$ and $y$ should both be a factor of $720$.

 After experimenting, we see this is achieved when $y=16$ and $x=15$, which means our value of $b$ is $720 \cdot \frac{16}{15} = 768$. Therefore, the sum is $7+6+8=\boxed{\textbf{(E) } 21}$.

 



\textbf{Solution 3 (Similar to previous solutions)}

To minimize the value of $b$, where it has to be an integer, and it has to be greater than 720, we can express the common ratio as $\dfrac{n+1}{n}$, where the value has to be greater than $1$, and $n$, and $n+1$ have to be factors of 720. Since the bigger the denominator gets, the smaller the value of the fraction, we essentially have to find the biggest value for $n$, where itself and $n+1$ are factors of $720$. From here, we can check whether $n(n+1) = 720$ yields an integer root, which it doesn't. So, then we check the next biggest factor of $720$, which is $360$. $n(n+1) = 360$, this doesn't have an integer root either. So, then we check the next biggest factor which is $240$, $n(n+1) = 240$, which we get $15$ as a root. This means the common ration is $\dfrac{16}{15}$. We then multiply $\dfrac{16}{15}$ times $720$ and add up the digits getting $\boxed{\textbf{(E) } 21}$.

 



\textbf{Solution 4}

Let the common ratio of the geometric sequence be $\frac{x}{y}$, with $x>y$. This means that $720(\frac{x}{y})$ and $720(\frac{y}{x})$ must both be integers, therefore $x$ and $y$ are both factors of $720$. We would achieve the smallest ratio $\frac{x}{y}$ if $x$ and $y$ are consecutive, so by listing out the factors of $720$, we find that $1$, $2$, $3$, $4$, $5$, $6$, $8$, $9$, $10$, $15$, $16$ are the only consecutive factors (any factors larger than these would result in the ratio simplifying, making it larger than just using consecutive integers). $15$ and $16$ are the largest, so we find the common ratio to be $\frac{16}{15}$, making $b=720(\frac{16}{15})$ giving us $768$. The sum of its digits is $\boxed{\textbf{(E) } 21}$.

 



\textbf{Solution 5 (similar to other solutions)}

We look for the smallest possible increasing ratio in the geometric sequence for $a$, $720$, and $b$.First, we find the prime factorization of $720$, which is really easy $720=2^4\cdot3^2\cdot5$. Then, we look for the closest possible fraction-ratio out of the primes, and after a quick look, it's easy to see it is $16$ and $15$. Now that we found our ratio, we multiply $720$ by $\frac{16}{15}$ to get $768$, so the sum of the digits of the least possible value of $b$ is $21$.

 NOTE: It's not 9/8 because we want it to be as close as possible to 1 ~gabrielle.esther

 


	\item (\textit{AMC 10 10A 2024}) Let $S$ be a subset of $\{1, 2, 3, \dots, 2024\}$ such that the following two conditions hold:


 
                \begin{itemize}
                    \item If $x$ and $y$ are distinct elements of $S$, then $|x-y| > 2.$
                \end{itemize}


                
                \begin{itemize}
                    \item If $x$ and $y$ are distinct odd elements of $S$, then $|x-y| > 6.$
                \end{itemize}


                What is the maximum possible number of elements in $S$?


 $\textbf{(A) }436 \qquad \textbf{(B) }506 \qquad \textbf{(C) }608 \qquad \textbf{(D) }654 \qquad \textbf{(E) }675$


 
 \textbf{Answer: }C



\textbf{Solution 1}

All lists are organized in ascending order:

 By listing out the smallest possible elements of subset $S,$ we can find that subset $S$ starts with $\{1, 4, 8, 11, 14, 18, 21, 24, 28, 31, \dots\}.$ It is easily noticed that the elements of the subset "loop around" every 3 elements, specifically adding 10 each time. This means that there will be $2024/10$ or $202R4$ whole loops in the subset $S,$ implying that there will be $202*3 = 606$ elements in S. However, we have undercounted, as we did not count the remainder that resulted from $2024/10$$.$ With a remainder of $4,$ we can fit $2$ more elements into the subset $S,$ namely $2021$ and $2024,$ resulting in a total of $606+2$ or $\boxed{\textbf{(C) }608}$ elements in subset $S$.

 
\textbf{NOTE:}

 To prove that this is the best we can do, consider adding each element one by one, for the first element, say n. If $n>2$, we can choose $n - 2$ which is always better. Therefore, $n = 1$ or $n = 2.$ 

 If $n = 2$ was optimal, then choose it, then the set of usable numbers in $S$ becomes 5 through 2024. We can transform the usable set of $S$ to $Q$ where $Q$ contains the numbers 1 through 2020. Because we assumed $n = 2$ was optimal, we can choose $n = 2$ for the set $Q$ too. Because every element in $Q$ is 4 below the elements of $S$, choosing 2 in $Q$ would mean choosing 6 in set $S$. By induction we see that our list would be $\{2,6,10,14,18,.....2022\}$ which only gives $506$ elements which is sub-optimal. Therefore, we can conclude that $n = 1$ is optimal, and we proceed as the solution above.

 -weihou0

 - minor latex edits by eqb5000

 



\textbf{Solution 2}

Notice that we can first place odd numbers, then place even numbers between each pair. We can start at $1$ and continue from there. Realize that the smallest number $k$ such that $kx+1$ reproduces odd number is $8$. The next ones are $10, 12, 14$. We can proceed to find the number of numbers in this particular sequence. From the equation $8x+1=2023$, we get that $x \leq 252.875$ works, so this means there is 253 solutions. Looking at $1,2,3,4,5,6,7,9$ we can see that there could only be 1 possible number between each pair, yielding $252+253=505$. Then see that we can fit two more into the number count since the set $2017$ to $2024$ can fit two evens. Now this means $A$ and $B$ don’t work. Now test out $10x+1$. Using the same method, we get that $608$ is the maximum number in the set. Everything above $x=10$ doesn’t work, as we can split it down into smaller subgroups, so the answer is $\boxed{\textbf{(C) }608}$.

 



\textbf{Solution 3}

We find the following pairs of numbers work: 

 
 \[(1,4), (4,8), (8,11), (11,14), (14, 18), (18, 21), (21, 24) \dots\]

 
Call a number an OE (odd-even) pair if the first number in the pair is odd and the second is even. Do the same for EE (even-even) and EO (even-odd) pairs. Notice that we're adding 10 to the pairs in each set, so there are 404 numbers in each set of pairs \textbf{excluding 1 and 2024:}

 
\textbf{OE pairs:} 
 \[(1,4), (11,14), ... , (2021,2024) \implies (202 + 1) \cdot 2 - 2 = 404\] 

 
\textbf{EE pairs:} 
 \[(4,8), (14,18), ..., (2014,2018)  \implies (201 + 1) \cdot 2 = 404\] 

 
\textbf{EO pairs:} 
 \[(8,11), (18,21), ..., (2018, 2021) \implies (201 + 1) \cdot 2 = 404\] 

 
Every number besides $1$ and $2024$ is being overcounted twice (for example, $4$ is counted twice in the EE and OE pairs), so we have $404 \cdot \frac{3}{2} = 606.$ Finally, we add $2$ (adding back $1$ and $2024$) to get $\boxed{\textbf{(C) }608}$.

 


	\item (\textit{AMC 10 10A 2024}) The numbers, in order, of each row and the numbers, in order, of each column of a $5 \times 5$ array of integers form an arithmetic progression of length $5$. The numbers in positions $(5, 5)$, $(2, 4)$, $(4, 3)$ and $(3, 1)$ are $0$, $48$, $16$, and $12$, respectively. What number is in position $(1, 2)$?

 \[\begin{bmatrix} . & ? &.&.&. \\ .&.&.&48&.\\ 12&.&.&.&.\\ .&.&16&.&.\\ .&.&.&.&0\end{bmatrix}\]


 

$\textbf{(A)}~19\qquad\textbf{(B)}~24\qquad\textbf{(C)}~29\qquad\textbf{(D)}~34\qquad\textbf{(E)}~39$


 
 \textbf{Answer: }C



\textbf{Solution 1 (Arithmetic Sequences)}

Start from the $0$ square. Going up, let the common difference be $a$, and going left, let the common difference be $b$. Therefore, we have 
 \[\begin{bmatrix} . & ? &.&.&4a \\ .&.&.&48&3a\\ 12&.&.&.&2a\\ .&.&16&.&a\\ 4b&3b&2b&b&0\end{bmatrix}\] Looking at the third column, we can see that the common difference going up is $16-2b$. We fill this in: 
 \[\begin{bmatrix} . & ? &64-6b&.&4a \\ .&.&48-4b&48&3a\\ 12&.&32-2b&.&2a\\ .&.&16&.&a\\ 4b&3b&2b&b&0\end{bmatrix}\] Looking at the second row, $48$ has two values beside it, so we can write 
 \[48=\dfrac{48-4b+3a}{2}\rightarrow96=48-4b+3a\rightarrow48=-4b+3a,\] and we can do the same with the third row, which gives 
 \[32-2b=\dfrac{12+2a}{2}\rightarrow32-2b=6+a\rightarrow26=a+2b.\] Now we have the system of equations 
 \[48=-4b+3a\]
 \[26=a+2b,\] and solving it gives $a=20,b=3$, therefore we can now fill in the grid with actual numbers. But before doing that, note that we're only looking for a value in the first row, and because we already have two known values in that row, we can find the common difference for that row and not focus on anything else. 

 Focusing only on the first row yields 
 \[\begin{bmatrix} . & ? &46&.&80\end{bmatrix}\] This means that the common difference from right to left is $\dfrac{80-46}{2}=17$. Therefore, the desired value is $46-17=\boxed{\text{(C) }29}$ ~Tacos\_are\_yummy\_1

 



\textbf{Solution 2 (Some Basic Algebra and Answer Choices)}

Assume the number in position $(3, 3)$ is $x$. 
 \[\begin{bmatrix} .&?&.&.&. \\ .&.&.&48&.\\ 12&.&x&.&.\\ .&.&16&.&.\\ .&.&.&.&0\end{bmatrix}\] Then, the integer in position $(2, 3)$ will be $2x-16$, since $2x-16$ and $16$ average out to x. (They form an arithmetic sequence).
 \[\begin{bmatrix} .&?&.&.&. \\ .&.&2x-16&48&.\\ 12&.&x&.&.\\ .&.&16&.&.\\ .&.&.&.&0\end{bmatrix}\]Similarly, the integer in position $(3, 2)$ is $0.5x+6$, and the integer in position $(2, 2)$ is $4x-80$. 
 \[\begin{bmatrix} .&?&.&.&. \\ .&4x-80&2x-16&48&.\\ 12&0.5x+6&x&.&.\\ .&.&16&.&.\\ .&.&.&.&0\end{bmatrix}\]This means the number $7.5x-166$ is in position (1, 2).
 \[\begin{bmatrix} .&7.5x-166&.&.&. \\ .&4x-80&2x-16&48&.\\ 12&0.5x+6&x&.&.\\ .&.&16&.&.\\ .&.&.&.&0\end{bmatrix}\]

 The only answer choice that makes $x$ an integer is $\boxed{\textbf{(C) }29}$

 



\textbf{Solution 3 (Bashy)}

Let's look at the bottom-left corner. There can be $3$ possible cases for the bottom row. Either the numbers are negative, all $0$, or positive. 

 Let's assume all the numbers are negative. We know that the numbers in between the $?$ will be large, so we want to minimize it, because the largest answer choice is only $39$.

 
 \[\begin{bmatrix} .&?&.&.&. \\ .&.&.&48&.\\ 12&.&.&.&.\\ .&.&16&.&.\\ -4&-3&-2&-1&0\end{bmatrix}\]

 To fill in the $1$st and $3$rd column, we simply find the common differences. We get:

 
 \[\begin{bmatrix} 28&?&70&.&. \\ 20&.&52&48&.\\ 12&.&34&.&.\\ 4&.&16&.&.\\ -4&-3&-2&-1&0\end{bmatrix}\]

 Finding the $?$ , we notice that it is $49$, which is bigger than any of the other answer choices, so the bottom row can't be negative. (As the common difference for the bottom row gets larger, so does the common difference for the $1$st and $3$rd column, making the $?$ bigger, and because the smallest answer is $39$ while the minimum value for negative numbers is $49$, it can't be negative.)

 Now let's assume that all the numbers are 0. By just filling in the $1$st and $3$rd column, we find that it is:

 
 \[\begin{bmatrix} 24 & ? &64&.&. \\ 18&.&48&48&.\\ 12&.&32&.&.\\ 6&.&16&.&.\\ 0&0&0&0&0\end{bmatrix}\]

 which makes the $?=44$. $44$ is not an answer choice, so the bottom row cannot be all $0$'s.

 Now, we know that the bottom row must be positive. Here's where the bashing begins. Lucky for us, the numbers must all be integers, cutting down on the number of cases we need to test. Starting with the bottom row having a common difference of $1$, we get

 
 \[\begin{bmatrix} . & ? &.&.&. \\ 16&30&44&48&.\\ 12&21&30&39&48\\ 8&12&16&20&24\\ 4&3&2&1&0\end{bmatrix}\]

 At this point we notice that the second row does not form an arithmetic progression, so we test out the next case, with the common difference being $2$. We get

 
 \[\begin{bmatrix} . & ? &.&.&. \\ 14&27&40&48&.\\ 12&20&28&36&44\\ 10&13&16&19&22\\ 8&6&4&2&0\end{bmatrix}\]

 At this point we notice that the second row does not form an arithmetic progression, so we test out the next case, with the common difference being $3$. We get

 
 \[\begin{bmatrix} 12 & 29 &46&63&80 \\ 12&24&36&48&60\\ 12&19&26&33&40\\ 12&14&16&18&20\\ 12&9&6&3&0\end{bmatrix}\]

 , giving us the answer $\boxed{\textbf{(C) }29}$

 



\textbf{Solution 4 (Arithmetic Sequences)}

Let the bottom row of the grid have arithmetic difference $d$. Then the matrix becomes 
 \[\begin{bmatrix} . & ? &.&.&4d \\ .&.&.&48&3d\\ 12&.&.&.&2d\\ .&.&16&.&d\\ .&.&.&.&0\end{bmatrix}\]Using the middle column, we have the middle item equal to $32-2d$.
 \[\begin{bmatrix} . & ? &.&.&4d \\ .&.&.&48&3d\\ 12&.&32-2d&.&2d\\ .&.&16&.&d\\ .&.&.&.&0\end{bmatrix}\]Filling in the middle row (which has a common difference of $10-d$)
 \[\begin{bmatrix} . & ? &.&.&. \\ .&.&.&48&.\\ 12&22-d&32-2d&42-3d&52-4d\\ .&.&16&.&.\\ 4d&3d&2d&d&0\end{bmatrix}\]The arithmetic sequence in the 4th column (2nd from the right) then has a common difference of $-3d-6$, so we
 \[\begin{bmatrix} . & ? &.&.&. \\ .&.&.&48&.\\ 12&22-d&32-2d&42-3d&52-4d\\ .&.&16&36-6d&.\\ 4d&3d&2d&30-9d&0\end{bmatrix} = \begin{bmatrix} . & ? &.&.&. \\ .&.&.&48&.\\ 12&22-d&32-2d&42-3d&52-4d\\ .&.&16&36-6d&.\\ 4d&3d&2d&d&0\end{bmatrix}\]From the cell in the 4th column, 5th row we have $30-9d=d$, or $d=3$. Finally, using the 2nd column, $\text{?} + 3d = 2(22-d)$, or $\text{?}+9=38 \Rightarrow \text{?}=\boxed{\text{(C) }29}$

 



\textbf{Solution 5}

We subtract $12$ from each element of the array and denote $a_{i,j}$ the number in position $(i,j), a_{3,1} = 0, a_{4,3} = 4,  a_{2,4} = 36,  a_{5,5} = - 12.$Denoting $a_{3,2} = x \implies$ that in third row we get $(0, x, 2x, 3x, 4x).$ We use the property of an arithmetic progression and get
 \[a_{4,4} = 2a_{3,4} - a_{2,4} = 6x - 36,\] 
 \[a_{4,5} = 2 a_{4,4}- a_{3,4} = 12x - 76,\]
 \[a_{5,5} = 2 a_{4,5}- a_{3,5} \implies  - 12 = 2(2x - 76) - 4x \implies x = 7 = a_{3,2},  a_{3,3}=14.\]
 \[\begin{bmatrix} .&.&.&.&. \\ .&.&.&36&.\\ 0&x&2x&3x&4x\\ .&.&4&6x-36&12x-76\\ .&.&.&.&-12\end{bmatrix}\]
 \[a_{2,3} = 2 a_{3,3}- a_{4,4} = 2 \cdot 14 - 4 = 24,\]
 \[a_{2,2} = 2 a_{2,3}- a_{2,4} = 2 \cdot 24 - 36 = 12,\]
 \[a_{1,2} = 2 a_{2,2}- a_{3,2} = 2 \cdot 12 - 7 = 17.\]

 
 \[\begin{bmatrix} .&17&.&&.& \\ .&12&24&36&.\\ 0&7&14&.&.\\ .&.&.&.&.\\ .&.&.&.&.\end{bmatrix}\]We add $12$ and get the number in position $(1,2)$ as $17 + 12 = 29.$

 \textbf{vladimir.shelomovskii@gmail.com, vvsss}

 



\textbf{Solution 6 (Bilinear Form with Coordinate Translation)}

A fundamental structural property of grids with arithmetic progressions in both rows and columns is that every entry can be expressed in the bilinear form:
 \[A_{m,n} = amn + bm + cn + d\]where $m$ and $n$ are the row and column indices. This representation captures the essential two-dimensional linear structure of the problem.

 \textbf{Solving the Problem} To minimize computational complexity, we employ coordinate translation. Set the known value $12$ at position $(3,1)$ as our origin by defining $x = m - 3$ and $y = n - 1$. The given data points become:

 
                \begin{itemize}
                    \item $A_{0,0} = 12$ (from position $(3,1)$)\item $A_{-1,3} = 48$ (from position $(2,4)$)\item $A_{1,2} = 16$ (from position $(4,3)$)\item $A_{2,4} = 0$ (from position $(5,5)$)
                \end{itemize}


                We seek the value at $(-2, 1)$, corresponding to the original position $(1,2)$.

 Substituting into $A_{x,y} = axy + bx + cy + d$:

 
            \begin{enumerate}
                \item $A_{0,0} = d = 12$\item $A_{-1,3} = -3a - b + 3c + 12 = 48 \implies -3a - b + 3c = 36$\item $A_{1,2} = 2a + b + 2c + 12 = 16 \implies 2a + b + 2c = 4$\item $A_{2,4} = 8a + 2b + 4c + 12 = 0 \implies 4a + b + 2c = -6$
            \end{enumerate}
            Subtracting equation (3) from equation (4) eliminates both $b$ and $c$:
 \[(4a + b + 2c) - (2a + b + 2c) = -6 - 4 \implies 2a = -10 \implies a = -5\]

 With $a = -5$ in equation (3): $-10 + b + 2c = 4 \implies b + 2c = 14$

 With $a = -5$ in equation (2): $15 - b + 3c = 36 \implies -b + 3c = 21$

 Adding these equations: $5c = 35 \implies c = 7$, hence $b = 0$.

 Our formula is therefore $A_{x,y} = -5xy + 7y + 12$. Evaluating at $(-2, 1)$:
 \[A_{-2,1} = -5(-2)(1) + 7(1) + 12 = 10 + 7 + 12 = 29\]

 The answer is $\boxed{\textbf{(C) } 29}$.

 \textbf{Verification of the Formula} To confirm this bilinear form universally satisfies the arithmetic progression property, let $A_{m,n} = amn + bm + cn + d$.

 For fixed row $m$, the column difference is:
 \[A_{m,n+1} - A_{m,n} = am(n+1) + bm + c(n+1) + d - (amn + bm + cn + d) = am + c\]This is constant with respect to $n$, confirming each row forms an arithmetic progression.

 For fixed column $n$, the row difference is:
 \[A_{m+1,n} - A_{m,n} = a(m+1)n + b(m+1) + cn + d - (amn + bm + cn + d) = an + b\]This is constant with respect to $m$, confirming each column forms an arithmetic progression.

 \textbf{Two Distinct Derivations of the Formula} The bilinear form can be derived through two fundamentally different mathematical perspectives, each revealing unique structural insights.

 \textbf{Derivation 1: Algebraic and Functional Perspective}

 This approach exploits the nested linear structure without invoking difference operators.

 Since every row is an arithmetic progression, for fixed $n$, the function $A_{m,n}$ must be linear in $m$:
 \[A_{m,n} = u_n \cdot m + v_n\]where coefficients $u_n$ and $v_n$ depend only on $n$.

 To determine the functional forms of $u_n$ and $v_n$, evaluate at specific rows:

 Setting $m = 0$: $A_{0,n} = v_n$. Since column $0$ is an arithmetic progression, $v_n$ must be linear in $n$. Write $v_n = cn + d$.

 Setting $m = 1$: $A_{1,n} = u_n + v_n$. Since column $1$ is also an arithmetic progression, $u_n + v_n$ must be linear in $n$. Given that $v_n$ is already linear, $u_n$ must itself be linear in $n$. Write $u_n = an + b$.

 Substituting back:
 \[A_{m,n} = (an + b)m + (cn + d) = amn + bm + cn + d\]

 This derivation reveals that the bilinear form emerges naturally from the compositional structure of linear dependencies in both dimensions.

 \textbf{Derivation 2: Discrete Calculus Perspective}

 This approach uses difference operators, illuminating the discrete analog of Clairaut's theorem on equality of mixed partial derivatives.

 Define forward difference operators:

 
                \begin{itemize}
                    \item $\Delta_m A_{m,n} = A_{m+1,n} - A_{m,n}$ (row difference)\item $\Delta_n A_{m,n} = A_{m,n+1} - A_{m,n}$ (column difference)
                \end{itemize}


                The arithmetic progression conditions state:

 
                \begin{itemize}
                    \item $\Delta_n A_{m,n} = r_m$ depends only on $m$\item $\Delta_m A_{m,n} = c_n$ depends only on $n$
                \end{itemize}


                Consider the mixed second difference $M_{m,n} = \Delta_m \Delta_n A_{m,n}$:

 
                \begin{itemize}
                    \item Applying $\Delta_m$ to $r_m$: $M_{m,n} = r_{m+1} - r_m$ depends only on $m$\item Applying $\Delta_n$ to $c_n$: $M_{m,n} = c_{n+1} - c_n$ depends only on $n$
                \end{itemize}


                Since $M_{m,n}$ depends only on $m$ and only on $n$ simultaneously, it must be constant. Denote $M_{m,n} = a$.

 Therefore:

 
                \begin{itemize}
                    \item $c_{n+1} - c_n = a \implies c_n = an + b$ (column differences form an AP)\item $r_{m+1} - r_m = a \implies r_m = am + c$ (row differences form an AP)
                \end{itemize}


                Reconstructing $A_{m,n}$ from $A_{0,0} = d$ via discrete path integration:
 \[A_{m,n} = A_{0,0} + \sum_{j=0}^{n-1} \Delta_n A_{0,j} + \sum_{i=0}^{m-1} \Delta_m A_{i,n}\]
 \[A_{m,n} = d + \sum_{j=0}^{n-1} r_0 + \sum_{i=0}^{m-1} c_n = d + nc + m(an + b) = amn + bm + cn + d\]

 This derivation establishes a rigorous framework connecting the commutativity of discrete differences to the bilinear structure, providing a powerful generalization technique for higher-dimensional problems.

 


 \textbf{Complex systems of equations} From the given information, we can ensure that we find the correct answer by creating a $4$-variable, $4$-equation system to solve for $(1,2)$.

 Let the common difference moving side to side from row $3$ be $d$, the common difference moving vertically from $48$ be $y$, the common difference moving vertically from $16$ be $z$, and the common difference moving side to side from $0$ be $j$.

 We can then write $4$ equations and solve for each difference.

 $12+3x=48+y$$16-z=12+2x$$16+z=0-2j$$48+3y=0-j$

 We can solve each variable in terms of $y$ procedurally until we can get $x$ in terms of $y$, and solve each from there. We get that $j=-3y-48$. We then plug this into equation $3$.

 Doing this, we find that $z=6y+80$.

 We then plug this into equation $2$ and find that $x=-3y-38$.

 We then substitute our value for $x$ into equation $1$. We find that $y=-15$.

 We then plug this value into our equations and see that $x=7$, $j=-3$, and $z=-10$.

 Now we can see that the number directly to the left of $48$ must be $36$, as we add $20$, or two times $z$. This means the common difference of this row is $12$. We move another digit to the left and see that the number directly above $(1,2)$ will be $24$. 

 The digit directly below this will be $19$, as we add $x$ to $12$ or $7$. This means that to get to $(1,2)$ we must add $5$ to $24$.

 $24+5=\boxed{29}$.

 


	\item (\textit{AMC 10 10A 2024}) Let $\mathcal K$ be the kite formed by joining two right triangles with legs $1$ and $\sqrt3$ along a common hypotenuse. Eight copies of $\mathcal K$ are used to form the polygon shown below. What is the area of triangle $\Delta ABC$? 
                    \begin{figure}[H]
                        \centering
                        \includegraphics[width=0.5\textwidth]{Problem/AMC_10_10A_2024_22_Problem_Figure_1.png}
                    \end{figure}
                    


 $\textbf{(A) }2+3\sqrt3\qquad\textbf{(B) }\dfrac92\sqrt3\qquad\textbf{(C) }\dfrac{10+8\sqrt3}{3}\qquad\textbf{(D) }8\qquad\textbf{(E) }5\sqrt3$


 
 \textbf{Answer: }B



\textbf{Solution 1}


                    \begin{figure}[H]
                        \centering
                        \includegraphics[width=0.5\textwidth]{Solution/AMC_10_10A_2024_22_Solution_Figure_1.png}
                    \end{figure}
                    

 First, we should find the length of $AB$. In order to do this, as we see in the diagram, it can be split into 4 equal sections. Since diagram $K$ shows us that it is made up of two ${30,60,90}$ triangles, then the triangle outlined in red must be a $30,60,90$ triangle, as ${30+30=60}$, and the two lines are perpendicular (it is proveable, but during competition, it is best to assume this is true, as the diagram is drawn pretty well to scale). Also, since we know the length of the longest side of the red triangle is ${\sqrt3}$, then the side we are looking for, which is outlined in blue, must be $\frac{3}{2}$ by the ${1,\sqrt3, 2}$ relationship of ${30,60,90}$ triangles. Therefore $AB$, which is the base of the triangle we are looking, for must be $6.$

 
Now all we have to do is find the height. We can split the height into 2 sections, the green and the light green. The green section must be ${\sqrt3}$, as $K$ shows us. Also, the light green section must be equal to ${\frac{\sqrt3}{2}}$, as in the previous paragraph, the triangle outlined in red is $30,60,90$. Then, the green section, which is the height, must be ${\sqrt3}+{\frac{\sqrt3}{2}}$, which is just ${\frac{3\sqrt3}{2}}$.

 
Then the area of the triangle must be ${\frac{1}{2}}\cdot {b} \cdot {h}$, which is just $\boxed{  \textbf{(B) } \frac{9}{2} \sqrt3}.$

 Simple Proof of the red triangle having $30^\circ$, $60^\circ$, and $90^\circ$ angles:

 In the small kite illustrated in the problem, we can see that the two angles on the left side add up to $60^\circ$ (this is easily noticeable if you look hard enough.) So, the angle that the arrow is pointing to is $60^\circ$. Since another angle in the red triangle is obviously $90^\circ$, the portion of the angle that is in the red triangle and at point A is $180^\circ - 90^\circ - 60^\circ = 30^\circ$. Therefore, the red triangle is a $30^\circ$ $60^\circ$ $90^\circ$ triangle.

 



\textbf{Solution 2}

Let $\mathcal K$ be quadrilateral $MNOP$. Drawing line $MO$ splits the triangle into $\Delta MNO$.Drawing the altitude from $N$ to point $Q$ on line $MO$, we know $NQ$ is $\frac{\sqrt{3}}{2}$, $MQ$ is $\frac{3}{2}$, and $QO$ is $\frac{1}{2}$.

 
                    \begin{figure}[H]
                        \centering
                        \includegraphics[width=0.5\textwidth]{Solution/AMC_10_10A_2024_22_Solution_Figure_2.png}
                    \end{figure}
                    

 Due to the many similarities present, we can find that $AB$ is $4(MQ)$, and the height of $\Delta ABC$ is $NQ+MN$

 $AB$ is $4(\frac{3}{2})=6$ and the height of $\Delta ABC$ is $\sqrt3+\frac{\sqrt{3}}{2}=\frac{3\sqrt{3}}{2}$.

 Solving for the area of $\Delta ABC$ gives $6\cdot\frac{3\sqrt{3}}{2}\cdot\frac{1}{2}$ which is $\boxed{\textbf{(B) }\dfrac92\sqrt3}$

 



\textbf{Solution 3}

Let's start by looking at kite $\mathcal K$. We can quickly deduce based off of the side lengths that the kite can be split into two $30-60-90$ triangles. Going back to the triangle $\triangle ABC$, focus on side $AB$. There are $4$ kites, they are all either reflected over the line $AB$ or a line perpendicular to $AB$, meaning the length of $AB$ can be split up into 4 equal parts. 

 Pick out the bottom-left kite, and we can observe that the kite and the triangle formed by the intersection of the kite and $\Delta ABC$ share a $60$ degree angle. (this was deduced from the $30-60-90$ triangles in the kite) The line AB and the right side of the kite are perpendicular, forming a $90^{\circ}$ angle. Because that is also a $30-60-90$ triangle with a hypotenuse of $\sqrt3$, so we find the length of AB to be $4*3/2$, which is $6$. 

 Then, we can drop an altitude from $C$ to $AB$. We know that will be equivalent to the sum of the longer side of the kite and the shorter side of the triangle formed by the intersection of the kite and $\Delta ABC$. (Look at the line formed on the left of $C$ that drops down to $AB$ if you are confused) We already have those values from the $30-60-90$ triangles, so we can just plug it into the triangle area formula, $bh/2$. We get 
 \[6\cdot\dfrac{\sqrt3+\frac{\sqrt3}{2}}{2}\rightarrow3\cdot(\sqrt3+\dfrac{\sqrt3}{2})\rightarrow3\cdot\dfrac{\sqrt3}{2}\rightarrow\boxed{\textbf{(B) }\dfrac92\sqrt3}\]

 



\textbf{Solution 4}


                    \begin{figure}[H]
                        \centering
                        \includegraphics[width=0.5\textwidth]{Solution/AMC_10_10A_2024_22_Solution_Figure_3.png}
                    \end{figure}
                    

 



\textbf{Solution 5}

Let the point of intersection of $AB$ and the kite with $A$ as vertex be $D$. 

 Let the left kite with $C$ as a vertex touch the kite with $A$ as vertex at point $E$. 

 $\triangle ADE$ is a $30-60-90$ so $AD = \frac{3}{2}$ and $DE = \frac{\sqrt3}{2}$.

 So, $AB = 4\cdot AD = 6$ and $CD=CE+DE= \frac{3\sqrt3}{2}$, and the area is $\frac12\cdot AB \cdot CD = \boxed{  \textbf{(B) } \frac{9}{2} \sqrt3}.$ 

 



\textbf{Solution 6(trig bash)}


                    \begin{figure}[H]
                        \centering
                        \includegraphics[width=0.5\textwidth]{Solution/AMC_10_10A_2024_22_Solution_Figure_4.png}
                    \end{figure}
                    
 \[\textbf{Step 1:}\]As stated in the solutions above we can easily find that $AB$ is split into $4$ equal parts, so we have $AB=4(AE)$ We can calculate $AE,$ by using $AA$ similarity to find $\triangle{ADE}$ is a $30-60-90$ triangle, therefore we have $AD=\sqrt{3},DE=\frac{\sqrt{3}}{2},$ and finally $AE=\frac{3}{2},$ therefore $AB=4(AE)=4(\frac{3}{2})=6.$ Similarly we have $CF$ is congruent to $FG,$ therefore $CF=FG=1.$ Next we have $AD$ is congruent to $DF$ telling us $AF=AD+DF=2(AD)=2\sqrt{3}.$ Noticing $\triangle{CFA}$ is right, we apply Pythagorean theorem to $\triangle{CFA}$ to find 
 \[AC^2=CF^2+AF^2\]
 \[AC^2=(1)^2+(2\sqrt{3})^2=13\]
 \[AC=\sqrt{13}.\]
 \[\textbf{Step 2:}\]Next we would like to calculate CJ, As said before $CF=FG=1$ so $CG=2.$ We know the inscribed angle between CG and GJ is $60\cdot2=120,$ and finally we know $GJ=1.$ So we apply LoC on triangle CGJ in order to find CJ. 
 \[CJ^2=CG^2+GJ^2-2\cdot{CG}\cdot{GJ}\cdot\cos{120^{\circ}}.\] 
 \[CJ^2=2^2+1^2-2\cdot2\cdot1\cdot\cos{120^{\circ}}.\]
 \[CJ^2=5-2\cdot2\cdot\frac{-1}{2}=7.\] 
 \[CJ=\sqrt{7}.\]
 \[\textbf{Step 3:}\]Now since we have all side lengths of $\triangle{CAJ}$ we can find $\cos{\angle{CAJ}}.$ Applying LoC again on $\triangle{CAJ}$ we have, 
 \[CJ^2=CA^2+AJ^2-2\cdot{CA}\cdot{AJ}\cdot\cos{\angle{CAJ}}.\] 
 \[{\sqrt{7}}^2={\sqrt{13}}^2+{3}^2-2\cdot3\cdot{\sqrt{13}}\cdot\cos{\angle{CAJ}}.\]
 \[7=22-6{\sqrt{13}}\cdot\cos{\angle{CAJ}}.\]
 \[\frac{-15}{-6{\sqrt{13}}}=\frac{5\sqrt{13}}{26}=\cos{\angle{CAJ}}.\]
 \[\textbf{Step 4:}\]We can solve for the area using the sin area formula which is $\frac{1}{2}\cdot{AC}\cdot{AB}\sin{\angle{CAJ}}.$ To find $\sin{\angle{CAJ}}$ we use the well known fact $\sin^2{x}+\cos^2{x}=1.$ So we find, 
 \[\sin^2{\angle{CAJ}}+\cos^2{\angle{CAJ}}=1.\]
 \[\sin^2{\angle{CAJ}}+(\frac{5\sqrt{13}}{26})^2=1.\]
 \[\sin^2{\angle{CAJ}}+\frac{25}{52}=1.\]
 \[\sin^2{\angle{CAJ}}=\frac{27}{52}.\]
 \[\sin{\angle{CAJ}}=\frac{3\sqrt{39}}{26}.\]Finally to wrap up we can find the area of $\triangle{ABC}$ using the sin area formula, 
 \[[ABC]=\frac{1}{2}\cdot{AC}\cdot{AB}\cdot\sin{\angle{CAJ}}.\]
 \[[ABC]=\frac{1}{2}\cdot\sqrt{13}\cdot{6}\cdot\frac{3\sqrt{39}}{26}=\frac{9\sqrt{3}}{2}.\]Therefore our answer is $\boxed{\textbf{(B) } \frac{9}{2} \sqrt3.}$

 



\textbf{Solution 7(coordinate bash)}

Assume we set A as $(0, 0)$. The line AB can be split into 4 equal sections, easily noticed in the picture. Not that the first equal section, starting at A. When we follow the line AB, we see that it's a right angle, noticing its a 30, 60, 90 triangle. Because the hypothenuse of that triangle is $\sqrt{3}$. Using the 30, 60, 90 ratio, we see that one of those equal sections is 3/2. Making the coordinate of B, $(6, 0)$. 

 Now, we must find the coordinate of C. Note that when we drop C down to the perpindicular of AB, there is a $x$ component to C, and a $y$ component to C. To find the X component, we see that it's one equal section found before as 3/2, and another section equal to 1, given in the question, making it $x$ = 5/2. The Y component can be seen as another two sections. The two green sections marked in solution 1. One is $\sqrt{3}/2$ due to the same 30 60 90 triangle we found in the first part of the solution. And the other part is $\sqrt{3}$, given in the question. Making $y$ = $\frac{3}{2} \sqrt3$. So we see that C = ($5/2$, $\frac{3}{2} \sqrt3$)

 Using shoelace theorem, we get to get  $\boxed{  \textbf{(B) } \frac{9}{2} \sqrt3}.$

 sorry for no diagrams D:

 
-marcus

 
Alternate way to calculate area: Notice once you get the y-coordinate of C, you can realize it's just the height of the triangle. So the area would be$\frac{1}{2}$ times base times height where the base is ${6}$ and the height is $\frac{3}{2}$ $\sqrt{3}$.~ilikemath247365

 \textbf{Sidenote} The figure made of kites was in fact, an einstein tile (also known as the Hat). The einstein tile was discovered in November of 2022 by a retired printing systems engineer in England. It comes from the German word for one stone, and has been proven to infinitely tile a plane aperiodically, quite a unique property - it took decades for mathematicians to find a periodical tilings of smaller and smaller numbers, with Penrose discovering a method using only 2 tilings in 1974.~LeonQS

 



\textbf{Solution 8}

When looking at the figure, we can see that the corner is shaped irregularly. This can make a regular hexagon with side length 2. Then, you would find the height of the triangle using the hexagon, which would get you \(\frac{1}{2} \times 6 \times \frac{3\sqrt{3}}{2}\), which is $\boxed{  \textbf{(B) } \frac{9}{2} \sqrt3}.$

 



\textbf{Solution 9}


                    \begin{figure}[H]
                        \centering
                        \includegraphics[width=0.5\textwidth]{Solution/AMC_10_10A_2024_22_Solution_Figure_5.png}
                    \end{figure}
                    

 The triangle with side lengths $1$, $\sqrt3$, and $2$ will be referenced as "half kite" in this solution

 First, notice how the altitude is composed of three half kites. 

 After finding the altitude of a half kite, which is $\frac{\sqrt3}{2}$, we can figure out that the altitude of triangle ABC is $\frac{3\sqrt3}{2}$

 Now we take a look at the base.Notice how we can create two half kites, each with base $2$. Adding that to the two $1$ unit segments and the base will have length 6. 

 Using the area formula for a triangle, we get $6\cdot\frac{3\sqrt{3}}{2}\cdot\frac{1}{2}$ which is $\boxed{\textbf{(B) }\dfrac92\sqrt3}$

 



\textbf{Solution 10 (Fast)}

Notice that $[ABC]=\frac{AB}{2}\cdot K$ where $K$ is the length of the altitude from $C$ to $\overline{AB}.$ Using the similar triangles shown below yields that $\frac{\overline{AB}}{2}=3$ and $K=\frac{3\sqrt{3}}{2}$ so $[ABC]=\boxed{\textbf{(B) }\frac 92\sqrt{3}}.$
                    \begin{figure}[H]
                        \centering
                        \includegraphics[width=0.5\textwidth]{Solution/AMC_10_10A_2024_22_Solution_Figure_6.png}
                    \end{figure}
                    ~shockfront99

 \textbf{Asymptote Diagram Reconstruction} I saw there was no asymptote code for this, so I tried my best to recreate the diagram. 
                    \begin{figure}[H]
                        \centering
                        \includegraphics[width=0.5\textwidth]{Solution/AMC_10_10A_2024_22_Solution_Figure_7.png}
                    \end{figure}
                    Feel free to use it :) (As you can see I used it to make the diagram in my solution)

 


	\item (\textit{AMC 10 10A 2024}) Integers $a$, $b$, and $c$ satisfy $ab + c = 100$, $bc + a = 87$, and $ca + b = 60$. What is $ab + bc + ca?$


 $\textbf{(A) }212 \qquad \textbf{(B) }247 \qquad \textbf{(C) }258 \qquad \textbf{(D) }276 \qquad \textbf{(E) }284 \qquad$


 
 \textbf{Answer: }D



\textbf{Solution 1}

Subtracting the first two equations yields $(a-c)(b-1)=13$. Notice that both factors are integers, so $b-1$ could equal one of $13,1,-1,-13$ and $b=14,2,0,-12$. We consider each case separately:

 For $b=0$, from the second equation, we see that $a=87$. Then $87c=60$, which is not possible as $c$ is an integer, so this case is invalid.

 For $b=2$, we have $2c+a=87$ and $ca=58$, which by experimentation on the factors of $58$ has no solution, so this is also invalid.

 For $b=14$, we have $14c+a=87$ and $ca=46$, which by experimentation on the factors of $46$ has no solution, so this is also invalid.

 Thus, we must have $b=-12$, so $a=12c+87$ and $ca=72$. Thus $c(12c+87)=72$, so $c(4c+29)=24$. We can simply trial and error this to find that $c=-8$ so then $a=-9$. The answer is then $(-9)(-12)+(-12)(-8)+(-8)(-9)=108+96+72=\boxed{\textbf{(D) }276}$.

 By Antony Huang ~Very Minor Edit by lucassf12

 



\textbf{Solution 2}

Adding up first two equations: 
 \[(a+c)(b+1)=187\]
 \[b+1=\pm 11,\pm 17\]
 \[b=-12,10,-18,16\]

 Subtracting equation 1 from equation 2: 
 \[(a-c)(b-1)=13\]
 \[b-1=\pm 1,\pm 13\]
 \[b=0,2,-12,14\]

 
 \[\Rightarrow b=-12\]

 Which implies that $a+c=-17$ from $(a+c)(b+1)=187$

 Giving us that $a+b+c=-29$

 Therefore, $ab+bc+ac=100+87+60-(a+b+c)=\boxed{\text{(D) }276}$

 



\textbf{Solution 3 (Guess and check)}

The idea is that you could guess values for $c$, since then $a$ and $b$ are factors of $100 - c$. The important thing to realize is that $a$, $b$, and $c$ are all negative. Then, this can be solved in a few minutes, giving the solution $(-9, -12, -8)$, which gives the answer $\boxed{\textbf{(D)} 276}$~andliu766

 


 



\textbf{Solution 4}


 \begin{align} ab + c &= 100 \\ bc + a &= 87 \\ ca + b &= 60 \end{align}

 
 \[(1) + (2) \implies  ab + c +bc + a = (a+c)(b+1)=187\implies b+1=\pm 11,\pm 17\]

 
 \[(1) - (2) \implies ab + c - bc - a = (a-c)(b-1)=13\implies b-1=\pm 1,\pm 13\]

 Note that $(b+1)-(b-1)=2$, and the only possible pair of results that yields this is $b-1=-13$ and $b+1=-11$, so $a+c=-17$. 

 Therefore, 

 



\textbf{Solution 5}


 \begin{align} ab + c &= 100 \\ bc + a &= 87 \\ ca + b &= 60 \end{align}

 \begin{align*}(1) - (2) \implies ab + c -bc - a &=(a-c)(b-1)=13 \\(2) - (3) \implies bc + a -ca - b &=(b-a)(c-1)=27 \\(3) - (1) \implies ca + b -ab - c &=(c-b)(a-1)=-40\end{align*}


 There are $3$ ordered triplets of $(a,b,c)$: $(5,14,4)$, $(-3,-12,-3)$, $(-9,-12,-8)$. 

 However, only the last ordered triplet meets all three equations.

 Therefore, $ab+ba+ac= -9\cdot-12+-12\cdot-8+-8\cdot-9 = \boxed{\textbf{(D) }276}.$

 



\textbf{Solution 6 (Elimination)}

Before we start, keep in mind that the problem is asking for the sum \(ab+bc+ac\). This is nothing but \(100+87+60-a-b-c\), or \(247-(a+b+c\)). 

 To solve the problem, we systematically test the options using elimination:

 Let's first check options A and B, since they only happen when a,b, and c sum to 35 or 0.We begin by testing three positive values, but none satisfy the equation when there is a plus sign. For example, \( (12, 8, 4) \) satisfies \( ab + c = 100 \), but does not satisfy \( bc + a = 87\), or \( ac + b = 60\). If \(a+b+c=0\), then not all of the numbers can be positive or negative, so this would not work.From this observation, we conclude that the answer cannot be \( \textbf{A} \) or \( \textbf{B} \).

 Now let's test the next option, option C.Option \( \textbf{C} \) states \( ab + bc + ca = 258 \). If true, then:

 \(a + b + c = -11\)

 This sum is too large. Furthermore, if all three numbers are negative, the solution still fails. For example, testing \( (-4, -5, -2) \) confirms the equation is not satisfied, as we get results that are too small. Thus, we eliminate option \( \textbf{C} \).

 Finally, let's test the last two options: D and E.For option \( \textbf{E} \), the sum \( a + b + c \) would be:

 \(247 - 284 = -37\)

 Testing values such as \( (-11, -12, -14) \), the resulting sums \( ab + c \), \( bc + a \), and \( ac + b \) are far too large to satisfy the equation.  Therefore, \( \textbf{E} \) is also eliminated.

 Once we have this answer, we still need to verify it by testing out numbers:Finally, we test option \( \textbf{D} \). Using \( ab + bc + ca = 276 \), we get that \(a+b+c = -29\). Also note that a, b, and c all have to be different, because the sums from the three equations are all different. We want to get the three closest values of a, b, and c such that they are all different, and the sum \(a+b+c = -29\). The values \( (-9, -12, -8) \) are the closest three numbers. When we try them, they satisfy all three equations.So, the correct answer is:$ab+ba+ac= -9\cdot-12+-12\cdot-8+-8\cdot-9 = \boxed{\textbf{(D) }276}.$

 


 \textbf{I NEED EASIER SOLUTIONS HERE:} 


	\item (\textit{AMC 10 10A 2024}) A bee is moving in three-dimensional space. A fair six-sided die with faces labeled $A^+, A^-, B^+, B^-, C^+,$ and $C^-$ is rolled. Suppose the bee occupies the point $(a,b,c).$ If the die shows $A^+$, then the bee moves to the point $(a+1,b,c)$ and if the die shows $A^-,$ then the bee moves to the point $(a-1,b,c).$ Analogous moves are made with the other four outcomes. Suppose the bee starts at the point $(0,0,0)$ and the die is rolled four times. What is the probability that the bee traverses four distinct edges of some unit cube?


 $\textbf{(A) }\frac{1}{54}\qquad\textbf{(B) }\frac{7}{54}\qquad\textbf{(C) }\frac{1}{6}\qquad\textbf{(D) }\frac{5}{18}\qquad\textbf{(E) }\frac{2}{5}$


 
 \textbf{Answer: }B



\textbf{Solution 1 (Simplest)}

We start by imagining the three dimensional plane.

 
                    \begin{figure}[H]
                        \centering
                        \includegraphics[width=0.5\textwidth]{Solution/AMC_10_10A_2024_24_Solution_Figure_1.png}
                    \end{figure}
                    Notice how the three dimensional plane is split into 8 different regions. There exists 8 cubes with one vertex on the point $(0,0,0)$. We need to consider each of these cases.

 We arbitrarily take a cube from one region of the three dimensional plane.

 Below is a sample drawing.

 
                    \begin{figure}[H]
                        \centering
                        \includegraphics[width=0.5\textwidth]{Solution/AMC_10_10A_2024_24_Solution_Figure_2.png}
                    \end{figure}
                    Roll 1: From $(0,0,0)$, there are 3 favorable outcomes that can occur. They are listed below.

 
                    \begin{figure}[H]
                        \centering
                        \includegraphics[width=0.5\textwidth]{Solution/AMC_10_10A_2024_24_Solution_Figure_3.png}
                    \end{figure}
                    Roll 2: WLOG, say the bee moved up one, there are then 2 ways the bee can go.

 
                    \begin{figure}[H]
                        \centering
                        \includegraphics[width=0.5\textwidth]{Solution/AMC_10_10A_2024_24_Solution_Figure_4.png}
                    \end{figure}
                    Roll 3: WLOG, say the bee moved forward one. From this point, there is again 2 ways the bee can go.

 
                    \begin{figure}[H]
                        \centering
                        \includegraphics[width=0.5\textwidth]{Solution/AMC_10_10A_2024_24_Solution_Figure_5.png}
                    \end{figure}
                    Roll 4: Finally, WLOG, say the bee moves to the left one. From that point, there are still 2 ways the bee can go.

 
                    \begin{figure}[H]
                        \centering
                        \includegraphics[width=0.5\textwidth]{Solution/AMC_10_10A_2024_24_Solution_Figure_6.png}
                    \end{figure}
                    The total probability for one cube is $\frac{3 \cdot 2 \cdot 2 \cdot 2}{6^4} \Rightarrow \frac{1}{54}$. We have 8 of these cubes, which gives us a total probability of $\frac{8}{54}$.

 However, we have overcounted. Notice that if the bee moved down one instead of left one at roll 3, it would only have one way. This is because the bee has a choice between either the origin or a unique point. This will always occur on the third move, and because we have 8 cubes(The 3d graph has one for each quadrant), the probability of this happening is $\frac{3 \cdot 8}{6^4} = \frac{1}{54}$*.

 Our probability is $\frac{8}{54} - \frac{1}{54} =$ $\boxed{\textbf{(B) }\frac{7}{54}}$.

 \textbf{*} One might question this logic here, onto why it is that we multiply the moves by the number of cubes over the number of rolls.

 The theory comes from expected value. It is certain that the bee will be selective in a situation where it only has one choice every 3 moves. This is due to the figure being a cube. Then, the expected probability of this happening is the moves multiplied by the single probability that this case appears, which is $\frac{1}{6^4}$. Finally, we multiply by 8, obtaining $\frac{3 \cdot 8}{6^4} = \frac{1}{54}$ and continuing as follows.

 



\textbf{Solution 2 (Higher End Mathematics)}

I put this as an experiment. If anything feels off or you don't understand, do not worry. This is very high end mathematics, and it is fine to feel confused. Solution 1 and the below solutions are much easier than this.

 Here, we explore the topological (study of bendable figures in space) definition of this problem.

 Consider an \( n \)-dimensional torus that the bee travels upon. Then, this torus must be an \( 8 \times 8 \times 8 \), as there are \( 8 \) three-dimensional cubes. The reason we can relate this to a torus is due to the physique of the \( 3D \) plane. 

 Then, we wish to find the ways to get from \( (0,0,0) \) to \( (0,0,0) \) in \( 512 \) moves. That isn't too difficult. In topology, the number of ways to do this is 

 
 \[\sum_{\gcd(x,y,z,8) = 1} \binom{8}{x} \binom{8}{y} \binom{8}{z}\]

 Summing over all values \( x, y, \) and \( z \) gives \( 2^{21} \times 7 \).

 Now, the total. The total number of moves that exist are all \( 512 \) plausible moves upon each probable space such that not all \( x, y, z \) are in \( \gcd \) relation of \( 1 \) with \( 8 \). There is then \( 3^{512} \) different ways to permute \( x, y, \) and \( z \). However, not all of these are plausible, as \( x + y + z = 512 \). These coordinates then force their way back to the origin due to the fact \( x \equiv y \equiv z \equiv 0 \pmod{8} \). Then, we can continue with roots of unity to filter out these multiples of \( 8 \). Our standard divisibility counter must satisfy 

 
 \[\mathbf{1}_{x \equiv 0 \pmod{8}} = \frac{1}{8} \sum_{j=0}^7 \omega^{jx}, \quad \omega = e^{2\pi i / 8}\]

 (Recall that this bolded 1 is the indicator function).

 Application to all coordinates satisfies 

 
 \[N = \frac{1}{8^3} \sum_{a,b,c=0}^{7} \sum_{x+y+z=512} \binom{512}{x, y, z} \omega^{ax+by+cz}\]

 Through the \href{https://artofproblemsolving.com/wiki/index.php?title=Multinomial\_Theorem}{Multinomial Theorem}, 

 
 \[\sum_{x+y+z=512} \binom{512}{x, y, z} \omega^{ax+by+cz} = (\omega^a + \omega^b + \omega^c)^{512}\]

 is simplified, and thus we obtain 

 
 \[N = \frac{1}{512} \sum_{a,b,c=0}^{7} (\omega^a + \omega^b + \omega^c)^{512}\]

 Because 512 is divisible by 8, any term where \( |\omega^a + \omega^b + \omega^c| < 2 \) doesn't contribute to the overall terms \( x \), \( y \), and \( z \). The only triplets that count are magnitude 2.

 Carrying this out gives \( 3^3 \) different full unital root contributions, and these are applied in \( 2^{22} \) different roots, for a total \( 2^{22} \times 27 \) ways. 

 Our probability is then
 \[\frac{2^{21} \times 7}{2^{22} \times 27} = \frac{7}{2 \times 27} = \boxed{\frac{7}{54}}\]

 Once again, this is an incredibly difficult solution. Before even attempting this, be sure you have a solid background in intermediate to Olympiad level algebra, precalculus, and counting/probability.

 



\textbf{Solution 3}

WLOG, assume that the first two moves are equal for all possible combinations, since the direction does not matter. The first move has a $\frac{6}{6}$ probability of being along one of the $8$ unit cubes around the origin, and the second move has a $\frac{4}{6}$ chance.Now, there are two cases. We are currently on one of the points of the $2$ by $2$ squares that are aligned with the axes.The first case is if the bee moves to the corner of a cube farthest away from the origin.Here, there is a $\frac{2}{6}$ chance of this happening and a $\frac{2}{6}$ chance of the fourth move remaining on one of the cubes.The second case is if the bee moves along the same plane of the $2$ by $2$ squares previously, ending up on a point 1 away from the origin. There is a $\frac{1}{6}$ chance of this happening and a $\frac{3}{6}$ chance of remaining on one of the cubes.Now, multiply and sum for the answer.
 \[\frac{2}{3}\cdot(\frac{1}{3}\cdot\frac{1}{3}+\frac{1}{6}\cdot\frac{1}{2})=\frac{2}{3}\cdot(\frac{1}{9}+\frac{1}{12})\]Evaluating this gives you the answer of $\boxed{\textbf{(B) }\frac{7}{54}}$.Solution by \href{https://artofproblemsolving.com/wiki/index.php?title=User:Juwushu}{juwushu}.Minor edits by andliu766

 



\textbf{Solution 4}

Remove the dice. It is not necessary for this problem. It is the same problem if you instead view this problem as the bee moving one up or down in the $x, y,$ and $z$ directions.

 Count the total number of ways for one cube of the 8 total ones. It is not hard to see that after the first step (which you can travel to $(1,0,0), (0,1,0),$ or $(0,0,1)$ -  3 ways), you have a total of 8 moves to traverse a face of that 1x1 cube. Hence, for each cube, you would have $8 \cdot 3 = 24$ total moves.

 There are a total of 8 cubes you could have traversed, though, so multiplying this by $8$ yields $8 \cdot 8 \cdot 3.$

 We have to account for overcounting. For example, when we traversed the face in the $xz$ plane $(0,0,0), (0,0,1),(1,0,0),(1,0,1),$ the number of ways that we traverse that particular face got counted twice, once from the cube in the far upper right and once from the cube in the closer upper right.

 For a given face, you can go clockwise or counterclockwise, hence yielding $2$ ways. There are a total of $12$ shared faces, each one shared by exactly two faces so no additonal PIE is needed. We simply subtract $2\cdot12$ from our original count now, yielding $8 \cdot 3 \cdot 8 - 2 \cdot 12 = 168.$ 

 With no restrictions, there are a total of $6^4$ ways to move four steps. Your probability is $\frac{168}{6^4}=\frac{7}{54}.$

 



\textbf{Solution 5}

An invariant is that the sum of all three coordinates change parity with every roll of the dice. Hence, to end up on a vertex of a unit cube, it can only be $(0,0,0)$, $(1,1,0)$, $(1,0,1)$, $(0,1,1)$, and all the negative permutations of each listed (in total, $13$ different final coordinates can be reached).

 For $(0,0,0)$, we choose two of the three coordinates to have 'cancelling' operations. Let these be $a$, $a'$, $b$, $b'$ (where $a$ cancels out $a'$, hence cannot be placed next to each other, which would be traversing the same edge). There are $3\cdot2\cdot2\cdot2=24$ ways for this (we can only have forms of $aba'b'$ or $bab'a'$ and their permutations among $a$ and $b$).

 For $(1,1,0)$, we must have two operations $a$, $b$ and a pair $c$, $c'$. Experimenting with having $3$ of the same type of operation leads to finding a repeated edge no matter the arrangement. These can be arranged in $12$ ways (place the operations $a$ and $b$ first, placing the operations $c$ and $c'$ in the remaining spaces).

 Analogously, $(-1,1,0)$, $(1,-1,0)$, $(-1,-1,0)$ all have the same number of ways, $12$ each (they are merely a rotation of all possibilities of $(1,1,0)$). Furthermore, we can generalise this to all the $8$ negative permutations of $(1,0,1)$, $(0,1,1)$, for a total of $12\cdot4\cdot3=144$. 

 So, $\frac{24+144}{6^4}=\frac{7\cdot24}{54\cdot24}=\frac{7}{54}$, hence we pick $\boxed{\textbf{(B) }\frac{7}{54}}$.

 -lisztepos

 


	\item (\textit{AMC 10 10A 2024}) The figure below shows a dotted grid $8$ cells wide and $3$ cells tall consisting of $1''\times1''$ squares. Carl places $1$-inch toothpicks along some of the sides of the squares to create a closed loop that does not intersect itself. The numbers in the cells indicate the number of sides of that square that are to be covered by toothpicks, and any number of toothpicks are allowed if no number is written. In how many ways can Carl place the toothpicks?


 
                    \begin{figure}[H]
                        \centering
                        \includegraphics[width=0.5\textwidth]{Problem/AMC_10_10A_2024_25_Problem_Figure_1.png}
                    \end{figure}
                    


 $\textbf{(A) }130\qquad\textbf{(B) }144\qquad\textbf{(C) }146\qquad\textbf{(D) }162\qquad\textbf{(E) }196$


 
 \textbf{Answer: }C



\textbf{Solution 1 (Best Motivated)}

Observations: 

 1. You can not have a vertical line in any place other than the first two columns and the last two columns. If we did, we would have at least one of the middle cells with toothpicks along more than one side, which would violate the conditions of the problem.

 2. There are two cases that look completely different. We can have a long horizontal box that spans all eight cells either on top of or below the middle cells, or we have to have a shape that looks like a rectangle, except with a few places "pushed" in. 

 Thus, using casework, we can split the task of finding those rectangles with squiggly edges into 3 cases.

 For case 1, we assume that the green lines shown below are given (always have toothpicks on them). In effect, we will use all eight columns.

 
                    \begin{figure}[H]
                        \centering
                        \includegraphics[width=0.5\textwidth]{Solution/AMC_10_10A_2024_25_Solution_Figure_1.png}
                    \end{figure}
                    

 
                    \begin{figure}[H]
                        \centering
                        \includegraphics[width=0.5\textwidth]{Solution/AMC_10_10A_2024_25_Solution_Figure_2.png}
                    \end{figure}
                    

 The only toothpicks we can place that will connect to the red lines are to go horizontally inward:

 
                    \begin{figure}[H]
                        \centering
                        \includegraphics[width=0.5\textwidth]{Solution/AMC_10_10A_2024_25_Solution_Figure_3.png}
                    \end{figure}
                    

 Now, concentrate on the first row of squares. A toothpick can be placed on either the bottom or top and connected to a continuous squiggle by adding vertical toothpicks:

 
                    \begin{figure}[H]
                        \centering
                        \includegraphics[width=0.5\textwidth]{Solution/AMC_10_10A_2024_25_Solution_Figure_4.png}
                    \end{figure}
                    

 How many squiggles are possible?

 
                    \begin{figure}[H]
                        \centering
                        \includegraphics[width=0.5\textwidth]{Solution/AMC_10_10A_2024_25_Solution_Figure_5.png}
                    \end{figure}
                    

 We can summarize this by giving a high squiggle position a 1 and a low position a 0, thus we have a 6-digit binary sequence. Thus, we can have $2^6=64$ ways to make this squiggle. (The binary is not absolutely necessary, but it works.)  

 Case 2: We can also pull in one of the sides, thus we can have a squiggle with 5 binary digits, which only uses the first or last 7 columns:

 
                    \begin{figure}[H]
                        \centering
                        \includegraphics[width=0.5\textwidth]{Solution/AMC_10_10A_2024_25_Solution_Figure_6.png}
                    \end{figure}
                    

 
                    \begin{figure}[H]
                        \centering
                        \includegraphics[width=0.5\textwidth]{Solution/AMC_10_10A_2024_25_Solution_Figure_7.png}
                    \end{figure}
                    

 Here, we only have 5 binary digits to work with, so there are $2^5=32$ ways to make this squiggle for each individual subcase. There are two subcases, one with the first 7 columns, and the other with the last 7, so we have a total of $32\cdot 2 = 64$ arrangements in this case.

 Case 3: We can use an even smaller section. Using only the middle 6 columns gives us a 4-wide squiggle:

 
                    \begin{figure}[H]
                        \centering
                        \includegraphics[width=0.5\textwidth]{Solution/AMC_10_10A_2024_25_Solution_Figure_8.png}
                    \end{figure}
                    

 Thus, there are $2^4=16$ ways to make this squiggle. 

 These three cases together cover all loops of this form. If we try to bring the square bracket like shapes on each side any closer, there will be some middle cells that do not touch any toothpicks. Adding up all our cases for these types of shapes: $64+32+32+16=144$.

 However, there are two more ways to draw a qualifying shape:

 
                    \begin{figure}[H]
                        \centering
                        \includegraphics[width=0.5\textwidth]{Solution/AMC_10_10A_2024_25_Solution_Figure_9.png}
                    \end{figure}
                    

 We can draw a rectangle like that in the first row or third row. Thus, we have a total of $144+2=\boxed{\textbf{(C) }146}$ ways.

 
A note to (potential) editors:This answer was not made to be concise or especially professional. It was made to explicitly explain this problem in a way so that it is easy to understand and follow.

 Notes:Remember these are the ONLY possible cases. It is impossible to cross through the rows of boxes of ones (in a snake like pattern) to connect the loop around the bottom since then the loop would intersect itself. 

 We are not undercounting by only counting the binary digits on the top. Consider the top squiggle. Each middle square must and can only come in contact with one toothpick, so if the top squiggle doesn't touch a specific middle cell, the bottom squiggle must to ensure that all squares touch one toothpick.~juwushu

 



\textbf{Solution 2 (Cheese)}

Notice that for any case where the closed loop does not connect from the top side of the ones and bottom side of the ones, there are two of these cases. A cheese solution can be found from this; noting that B and C are the only two options close to each other, and, being two apart, with people likely to forget this case, $\boxed{\textbf{(C) }146}$ is likely to be the correct answer.Cheese solution done by \href{https://artofproblemsolving.com/wiki/index.php?title=User:Juwushu}{juwushu}.

 



\textbf{Solution 3 (Observation)}

We have $2$ cases where the loop does not go through the middle. 

 If the loop goes through the middle, we must have a full column on $(0,8),(0,7),(1,8),(1,7).$ Then we have $6,5,5,4$ empty middle squares. For each one we can have one on top or one on bottom, so $2^6+2^5+2^5+2^4=144.$ Notice that for each case of fixed toothpicks, there is only one way to form the loop. Then we just add $2+144=\boxed{\textbf{(C) } 146.}$

 


\end{enumerate}